\documentclass[a4paper,12pt]{article}
\usepackage{cmap}
% --- Кодировки и шрифты ---
\usepackage[utf8]{inputenc}
\usepackage[T1, T2A]{fontenc} 
\usepackage[russian]{babel}

% --- Математические пакеты ---
\usepackage{amsmath}
\usepackage{amssymb}
\usepackage{amsthm}
\usepackage{mathtools}
\usepackage{tikz}
\usepackage{mathrsfs}

% --- Настройка полей ---
\usepackage[left=2cm, right=2cm, top=2cm, bottom=2cm]{geometry}

% --- Гиперссылки ---
\usepackage[colorlinks=true, linkcolor=blue, urlcolor=blue]{hyperref}

% --- ДОБАВЛЕННЫЕ ПАКЕТЫ ДЛЯ ОФОРМЛЕНИЯ ---
\usepackage[dvipsnames]{xcolor}
\usepackage[most]{tcolorbox}

% --- КОЛОНТИТУЛЫ (НАВИГАЦИЯ) ---
\usepackage{fancyhdr}
\pagestyle{fancy}
\fancyhf{}
\setlength{\headheight}{15pt}

% Команда для обновления кнопок
\newcommand{\setupnav}[2]{
	\fancyhead[L]{\hyperlink{toc}{\textbf{[Меню]}}}
	\fancyhead[R]{
		\ifx&#1& \textcolor{gray}{[$\leftarrow$ Назад]} \else \hyperlink{bilet#1}{[$\leftarrow$ Назад]} \fi
		\quad | \quad
		\ifx&#2& \textcolor{gray}{[Вперед $\rightarrow$]} \else \hyperlink{bilet#2}{[Вперед $\rightarrow$]} \fi
	}
}
\fancyfoot[C]{\thepage}

% --- ЦВЕТА ---
\definecolor{defmyorange}{RGB}{255, 140, 0}
\definecolor{thmblue}{RGB}{0, 102, 204}
\definecolor{proofpurp}{RGB}{128, 0, 128}
\definecolor{highlighter}{RGB}{255, 215, 0}
\definecolor{exgray}{RGB}{100, 100, 100}

% --- Определение окружений ---
\theoremstyle{plain}
\newtheorem{theorem}{Теорема}[section]
\newtheorem{lemma}[theorem]{Лемма}
\newtheorem{proposition}[theorem]{Утверждение}
\newtheorem{corollary}[theorem]{Следствие}

\theoremstyle{definition}
\newtheorem{definition}{Определение}[section]
\newtheorem{example}{Пример}[section]

\theoremstyle{remark}
\newtheorem{remark}{Замечание}[section]
\newtheorem*{idea}{Идея доказательства}

% --- МАГИЯ ПЕРЕОПРЕДЕЛЕНИЯ ---
\renewenvironment{definition}[1][]{
	\begin{tcolorbox}[enhanced, colback=defmyorange!5, colframe=defmyorange, coltitle=white, fonttitle=\bfseries, title={Определение: #1}, attach boxed title to top left={yshift=-2mm, xshift=2mm}, boxed title style={colback=defmyorange, sharp corners}, sharp corners, top=12pt, breakable]
	}{\end{tcolorbox}}

\renewenvironment{theorem}[1][]{
	\begin{tcolorbox}[enhanced, colback=thmblue!5, colframe=thmblue, fonttitle=\bfseries, title={Теорема #1}, sharp corners=south, drop shadow, breakable]
	}{\end{tcolorbox}}

\renewenvironment{lemma}[1][]{
	\begin{tcolorbox}[enhanced, colback=thmblue!5, colframe=thmblue, fonttitle=\bfseries, title={Лемма #1}, sharp corners=south, breakable]
	}{\end{tcolorbox}}

\renewenvironment{proposition}[1][]{
	\begin{tcolorbox}[enhanced, colback=thmblue!5, colframe=thmblue, fonttitle=\bfseries, title={Утверждение #1}, sharp corners=south, breakable]
	}{\end{tcolorbox}}

\renewenvironment{corollary}[1][]{
	\begin{tcolorbox}[enhanced, colback=thmblue!2, colframe=thmblue!70, fonttitle=\bfseries, title={Следствие #1}, sharp corners, breakable]
	}{\end{tcolorbox}}

\renewenvironment{proof}{
	\begin{tcolorbox}[
		enhanced, 
		breakable, 
		colback=proofpurp!2,    % Светлый фон
		colframe=proofpurp,     % Цвет рамки
		coltitle=white,         % Цвет текста в заголовке
		fonttitle=\bfseries, 
		title={Доказательство}, % Заголовок в рамке
		sharp corners,          % Острые углы как у теорем
		boxrule=1pt,            % Толщина рамки
		toptitle=2mm,           % Отступы в заголовке
		bottomtitle=2mm
		]
	}{\qed \end{tcolorbox}}

\renewenvironment{example}[1][]{
	\begin{tcolorbox}[enhanced, colback=exgray!5, colframe=exgray, fonttitle=\bfseries, title={Пример #1}, sharp corners, boxrule=0.5pt, breakable]
	}{\end{tcolorbox}}

\renewenvironment{remark}[1][]{
	\begin{tcolorbox}[
		enhanced,
		breakable,
		colback=exgray!4,
		colframe=exgray,
		fonttitle=\bfseries,
		title={Замечание #1},
		sharp corners,
		boxrule=0.5pt
		]
	}{
	\end{tcolorbox}
}

\setcounter{tocdepth}{2}


% --- НОВЫЙ ЦВЕТ ДЛЯ ИДЕИ ---
\definecolor{ideared}{RGB}{196, 30, 58} % Красивый красный (кардинал)

% --- МАГИЯ ПЕРЕОПРЕДЕЛЕНИЯ IDEA ---
\renewenvironment{idea}[1][]{
	\begin{tcolorbox}[
		enhanced,
		breakable,
		colback=ideared!5, 
		colframe=ideared, 
		borderline west={2pt}{0pt}{ideared},
		coltitle=ideared,
		fonttitle=\bfseries,
		title={Идея доказательства: #1},
		attach title to upper,
		after title={\par\nobreak\vspace{1mm}},
		sharp corners
		]
	}{\end{tcolorbox}}


% --- Заголовок ---
\title{Билеты по функциональному анализу--2}
\author{}
\date{\today}

\begin{document}
	
	\maketitle
	\hypertarget{toc}{} 
	
	\begin{abstract}
		Билеты по функциональному анализу для подготовки к экзамену. Оформлено по разделам программы (2026).
	\end{abstract}
	
	\tableofcontents
	\newpage
	
	\setcounter{section}{6}
	
	% VII. СЛАБАЯ СХОДИМОСТЬ В БАНАХОВОМ ПРОСТРАНСТВЕ
	\section{Слабая сходимость в банаховом пространстве}
	\hypertarget{bilet1}{}\setupnav{}{2} \providecommand{\weakconv}{\rightharpoonup}
\providecommand{\K}{\mathbb{K}}

\subsection{Изометричность вложения \(E\) в \(E^{**}\). Критерий слабой сходимости последовательности}

\begin{definition}[Слабая сходимость]
	Пусть \(E\) --- нормированное пространство. Последовательность
	\(\{x_n\}\subset E\) называется \textbf{слабо сходящейся} к \(x\in E\), если
	\[
	\forall f\in E^*
	\qquad
	f(x_n)\to f(x).
	\]
	Обозначение:
	\[
	x_n\weakconv x.
	\]
\end{definition}

\begin{definition}[Каноническое вложение \(E\) в \(E^{**}\)]
	Пусть \(E\) --- нормированное пространство. Для каждого \(x\in E\) определим функционал
	\[
	F_x:E^*\to\K,
	\qquad
	F_x(f)=f(x).
	\]
	Отображение
	\[
	J:E\to E^{**},
	\qquad
	Jx=F_x,
	\]
	называется \textbf{каноническим вложением} пространства \(E\) в \(E^{**}\).
\end{definition}

\begin{theorem}[Изометричность канонического вложения]
	Каноническое вложение
	\[
	J:E\to E^{**},
	\qquad
	(Jx)(f)=f(x),
	\]
	является линейной изометрией:
	\[
	\|Jx\|_{E^{**}}=\|x\|_E
	\qquad
	\forall x\in E.
	\]
\end{theorem}

\begin{proof}
	
	\textbf{1. Линейность отображения \(J\).}
	
	Для любых \(x,y\in E\), \(\alpha,\beta\in\K\) и \(f\in E^*\) имеем
	\[
	J(\alpha x+\beta y)(f)
	=
	f(\alpha x+\beta y)
	=
	\alpha f(x)+\beta f(y)
	=
	(\alpha Jx+\beta Jy)(f).
	\]
	Значит,
	\[
	J(\alpha x+\beta y)=\alpha Jx+\beta Jy.
	\]
	
	\textbf{2. Верхняя оценка нормы.}
	
	Для любого \(f\in E^*\) с \(\|f\|\leqslant 1\):
	\[
	|(Jx)(f)|
	=
	|f(x)|
	\leqslant
	\|f\|\,\|x\|
	\leqslant
	\|x\|.
	\]
	
	Следовательно,
	\[
	\|Jx\|
	=
	\sup_{\|f\|\leqslant 1}|(Jx)(f)|
	\leqslant
	\|x\|.
	\]о
	
	\textbf{3. Нижняя оценка нормы.}
	
	Если \(x=0\), то равенство очевидно. Пусть \(x\ne0\). По следствию из теоремы Хана--Банаха
	существует функционал \(f_0\in E^*\), такой что
	\[
	\|f_0\|=1,
	\qquad
	|f_0(x)|=\|x\|.
	\]
	Тогда
	\[
	\|Jx\|
	=
	\sup_{\|f\|\leqslant1}|(Jx)(f)|
	\geqslant
	|(Jx)(f_0)|
	=
	|f_0(x)|
	=
	\|x\|.
	\]
	
	Итак,
	\[
	\|Jx\|\leqslant \|x\|
	\qquad
	\text{и}
	\qquad
	\|Jx\|\geqslant \|x\|.
	\]
	Следовательно,
	\[
	\|Jx\|=\|x\|.
	\]
	Отображение \(J\) линейно и сохраняет норму, значит, \(J\) --- линейная изометрия.
\end{proof}

\begin{theorem}[Критерий слабой сходимости последовательности]
	Пусть \(E\) --- нормированное пространство.
	
	Последовательность \(x_n\) слабо сходится к \(x\) тогда и только
	тогда, когда выполнены два условия:
	
	\begin{enumerate}
		\item Последовательность норм ограничена:
		\(\exists M:\ \sup_n \|x_n\|\leqslant M.\)
		
		\item Сходимость есть на плотном множестве:
		\(\exists S\subset E^*,\ \overline{[S]}=E^*\) такое, что
		\(\forall s\in S:\ s(x_n)\to s(x).\)
	\end{enumerate}
\end{theorem}

\begin{proof}
	\textbf{1. Переход к операторам \(Jx_n\).}
	
	По определению слабой сходимости и канонического вложения \(J\):
	\[
	x_n\weakconv x
	\quad\Longleftrightarrow\quad
	\forall f\in E^*:\ 
	(Jx_n)(f)\to(Jx)(f).
	\]
	То есть
	\[
	x_n\weakconv x
	\quad\Longleftrightarrow\quad
	Jx_n\to Jx
	\quad\text{поточечно на }E^*.
	\]
	
	\textbf{2. Применение критерия поточечной сходимости.}
	
	Имеем
	\[
	Jx_n,Jx\in E^{**}=\mathcal{L}(E^*,\K).
	\]
	Пространство \(E^*\) банахово, а по условию
	\[
	\overline{[S]}=E^*.
	\]
	По критерию поточечной сходимости
	\[
	Jx_n\to Jx
	\]
	поточечно тогда и только тогда, когда
	\[
	\{\|Jx_n\|\}\ \text{ограничена}
	\]
	и
	\[
	(Jx_n)(s)\to(Jx)(s)
	\qquad
	\forall s\in S.
	\]
	
	\textbf{3. Возврат к \(x_n\).}
	
	Так как \(J\) --- изометрия и
	\[
	(Jx_n)(s)=s(x_n),
	\]
	получаем
	\[
	\{\|x_n\|\}\ \text{ограничена},
	\qquad
	s(x_n)\to s(x)
	\quad
	\forall s\in S.
	\]
	
	Следовательно,
	\[
	x_n\weakconv x
	\quad\Longleftrightarrow\quad
	\begin{cases}
		\{\|x_n\|\}\ \text{ограничена},\\
		s(x_n)\to s(x),\quad \forall s\in S.
	\end{cases}
	\]
\end{proof}\newpage
	\hypertarget{bilet2}{}\setupnav{1}{3} \providecommand{\weakconv}{\rightharpoonup}
\providecommand{\K}{\mathbb{K}}

\subsection{Слабая сходимость и линейные ограниченные операторы. Слабо ограниченные множества. Теорема Хана}

\begin{theorem}[Слабая сходимость и ограниченный линейный оператор]
	Пусть \(E_1,E_2\) --- нормированные пространства,
	\[
	A\in \mathcal{L}(E_1,E_2).
	\]
	Если
	\[
	x_n\weakconv x
	\qquad \text{в } E_1,
	\]
	то
	\[
	Ax_n\weakconv Ax
	\qquad \text{в } E_2.
	\]
\end{theorem}

\begin{proof}
	Докажем утверждение прямо по определению слабой сходимости.
	
	\textbf{1. Сведение к функционалу на \(E_1\).}
	
	Возьмём произвольный функционал
	\[
	g\in E_2^*.
	\]
	Чтобы доказать слабую сходимость
	\[
	Ax_n\weakconv Ax
	\quad \text{в } E_2,
	\]
	нужно проверить, что
	\[
	g(Ax_n)\to g(Ax).
	\]
	
	Рассмотрим композицию
	\[
	f=g\circ A,
	\qquad
	f(y)=g(Ay),
	\qquad
	y\in E_1.
	\]
	Тогда \(f\) линеен как композиция линейных отображений. Кроме того, для любого \(y\in E_1\)
	\[
	|f(y)|
	=
	|g(Ay)|
	\leqslant
	\|g\|\cdot\|Ay\|
	\leqslant
	\|g\|\cdot\|A\|\cdot\|y\|.
	\]
	Следовательно,
	\[
	f\in E_1^*.
	\]
	
	\textbf{2. Использование слабой сходимости в \(E_1\).}
	
	Так как
	\[
	x_n\weakconv x
	\qquad \text{в } E_1,
	\]
	то для функционала \(f=g\circ A\in E_1^*\) имеем
	\[
	f(x_n)\to f(x).
	\]
	То есть
	\[
	(g\circ A)(x_n)\to (g\circ A)(x),
	\]
	или, что то же самое,
	\[
	g(Ax_n)\to g(Ax).
	\]
	
	Функционал \(g\in E_2^*\) был произвольным, значит
	\[
	\forall g\in E_2^*
	\qquad
	g(Ax_n)\to g(Ax).
	\]
	По определению слабой сходимости это означает
	\[
	Ax_n\weakconv Ax.
	\]
\end{proof}

\begin{definition}[Слабо ограниченное множество]
	Пусть \(E\) --- нормированное пространство. Множество
	\[
	S\subset E
	\]
	называется \textbf{слабо ограниченным}, если для любого функционала
	\(f\in E^*\) множество
	\[
	f(S)=\{f(x):x\in S\}
	\]
	ограничено в \(\K\).
	
	Иными словами,
	\[
	\forall f\in E^*
	\quad
	\exists C(f)>0
	\quad
	\forall x\in S
	\qquad
	|f(x)|\leqslant C(f).
	\]
	Важно: константа \(C(f)\) может зависеть от функционала \(f\), но не зависит от точки \(x\in S\).
\end{definition}

\begin{remark}[Сравнение с обычной ограниченностью]
	Обычная ограниченность множества \(S\subset E\) означает:
	\[
	\exists C>0
	\quad
	\forall x\in S
	\qquad
	\|x\|\leqslant C.
	\]
	
	Слабая ограниченность означает:
	\[
	\forall f\in E^*
	\quad
	\exists C(f)>0
	\quad
	\forall x\in S
	\qquad
	|f(x)|\leqslant C(f).
	\]
	
	Теорема Хана утверждает, что эти два условия на самом деле эквивалентны.
\end{remark}

\begin{theorem}[Хана]
	Пусть \(E\) --- нормированное пространство и \(S\subset E\). Тогда
	\[
	S \text{ слабо ограничено}
	\quad\Longleftrightarrow\quad
	S \text{ ограничено}.
	\]
\end{theorem}

\begin{proof}
	Докажем обе импликации.
	
	\textbf{1. Ограниченность \(\Rightarrow\) слабая ограниченность.}
	
	Пусть \(S\) ограничено. Тогда существует \(R>0\), такое что
	\[
	\|x\|\leqslant R
	\qquad
	\forall x\in S.
	\]
	Возьмём произвольный функционал \(f\in E^*\). Тогда для любого \(x\in S\)
	\[
	|f(x)|
	\leqslant
	\|f\|\cdot\|x\|
	\leqslant
	\|f\|\cdot R.
	\]
	Значит, множество \(f(S)\) ограничено. Так как \(f\in E^*\) был произвольным,
	то \(S\) слабо ограничено.
	
	\textbf{2. Слабая ограниченность \(\Rightarrow\) ограниченность.}
	
	Докажем от противного. Пусть \(S\) слабо ограничено, но не ограничено.
	
	\begin{itemize}
		\item \textbf{Выберем быстро растущую последовательность.}
		
		Из неограниченности \(S\) следует, что для каждого \(n\in\mathbb{N}\)
		можно выбрать
		\[
		x_n\in S
		\]
		так, что
		\[
		\|x_n\|>n^2.
		\]
		Положим
		\[
		y_n=\frac{x_n}{n}.
		\]
		
		\item \textbf{Покажем, что \(y_n\weakconv 0\).}
		
		Возьмём произвольный функционал \(f\in E^*\). Так как \(S\) слабо ограничено,
		то существует \(C(f)>0\), такое что
		\[
		|f(x)|\leqslant C(f)
		\qquad
		\forall x\in S.
		\]
		В частности,
		\[
		|f(x_n)|\leqslant C(f).
		\]
		Тогда
		\[
		|f(y_n)|
		=
		\left|f\left(\frac{x_n}{n}\right)\right|
		=
		\frac{1}{n}|f(x_n)|
		\leqslant
		\frac{C(f)}{n}
		\to 0.
		\]
		Следовательно,
		\[
		f(y_n)\to 0=f(0)
		\qquad
		\forall f\in E^*.
		\]
		По определению слабой сходимости
		\[
		y_n\weakconv 0.
		\]
		
		\item \textbf{Получим противоречие с ограниченностью слабо сходящейся последовательности.}
		
		По лемме из предыдущего билета всякая слабо сходящаяся последовательность ограничена.
		Значит, последовательность \(\{y_n\}\) должна быть ограничена.
		
		Но по построению
		\[
		\|y_n\|
		=
		\left\|\frac{x_n}{n}\right\|
		=
		\frac{1}{n}\|x_n\|
		>
		\frac{1}{n}\cdot n^2
		=
		n.
		\]
		Следовательно,
		\[
		\|y_n\|>n
		\qquad
		\forall n\in\mathbb{N}.
		\]
		То есть \(\{y_n\}\) не ограничена. Противоречие.
	\end{itemize}
	
	Значит, предположение о неограниченности \(S\) неверно. Следовательно,
	\[
	S \text{ ограничено}.
	\]
\end{proof}

\newpage
	\hypertarget{bilet3}{}\setupnav{2}{4} \providecommand{\weakconv}{\rightharpoonup}
\providecommand{\K}{\mathbb{K}}

\subsection{Замкнутый шар в гильбертовом пространстве слабо секвенциально компактен (теорема Банаха)}

\begin{definition}[Слабая секвенциальная компактность]
	Пусть \(E\) --- нормированное пространство. Множество
	\[
	M\subset E
	\]
	называется \textbf{слабо секвенциально компактным}, если из любой последовательности
	\[
	\{x_n\}\subset M
	\]
	можно выделить подпоследовательность \(\{x_{n_k}\}\), слабо сходящуюся к некоторому
	элементу \(x\in M\):
	\[
	x_{n_k}\weakconv x.
	\]
\end{definition}

\begin{remark}[Слабая сходимость в гильбертовом пространстве]
	Пусть \(H\) --- гильбертово пространство. По теореме Рисса--Фреше каждый функционал
	\(f\in H^*\) имеет вид
	\[
	f(z)=(z,y)
	\]
	для некоторого \(y\in H\). Поэтому
	\[
	x_n\weakconv x
	\quad\Longleftrightarrow\quad
	\forall y\in H
	\qquad
	(x_n,y)\to(x,y).
	\]
\end{remark}

\begin{theorem}[Банаха]
	Пусть \(H\) --- гильбертово пространство и
	\[
	\overline{B}_R(0)=\{x\in H:\|x\|\leqslant R\}.
	\]
	Тогда замкнутый шар \(\overline{B}_R(0)\) слабо секвенциально компактен.
	
	Иными словами, из любой последовательности
	\[
	x_n\in\overline{B}_R(0)
	\]
	можно выделить подпоследовательность \(\{x_{n_k}\}\), такую что
	\[
	x_{n_k}\weakconv x
	\]
	для некоторого
	\[
	x\in\overline{B}_R(0).
	\]
\end{theorem}

\begin{idea}
	Берём последовательность в шаре. Она может жить в несепарабельном пространстве,
	но сама эта последовательность порождает сепарабельное подпространство
	\[
	L=\overline{\operatorname{span}}\{x_1,x_2,\ldots\}.
	\]
	Внутри \(L\) диагональным методом добиваемся сходимости всех чисел
	\[
	(x_m,x_{n_k})
	\]
	при фиксированном \(m\). Это даёт сходимость скалярных произведений на плотном
	подпространстве.
	
	После этого строим функционал
	\[
	F(z)=\lim_{k\to\infty}(z,x_{n_k}),
	\]
	по теореме Рисса--Фреше предъявляем элемент \(x\in L\), и получаем слабую сходимость.
\end{idea}

\begin{proof}
	Пусть
	\[
	x_n\in\overline{B}_R(0),
	\qquad n\in\mathbb{N}.
	\]
	То есть
	\[
	\|x_n\|\leqslant R
	\qquad
	\forall n.
	\]
	Нужно выделить подпоследовательность, слабо сходящуюся к элементу того же шара.
	
	\textbf{1. Переход к подпространству, порождённому последовательностью.}
	
	Рассмотрим
	\[
	L=\overline{\operatorname{span}}\{x_1,x_2,\ldots\}\subset H.
	\]
	Тогда
	\[
	x_n\in L
	\qquad
	\forall n.
	\]
	Подпространство \(L\) замкнуто в \(H\), поэтому \(L\) само является гильбертовым пространством.
	Кроме того, множество
	\[
	L_0=\operatorname{span}\{x_1,x_2,\ldots\}
	\]
	плотно в \(L\).
	
	Достаточно сначала построить слабую сходимость в \(L\), а затем перейти ко всему \(H\)
	через ортогональное разложение
	\[
	H=L\oplus L^\perp.
	\]
	
	\textbf{2. Диагональный выбор подпоследовательности.}
	
	Для фиксированного \(m\in\mathbb{N}\) числовая последовательность
	\[
	(x_m,x_n),
	\qquad n=1,2,\ldots,
	\]
	ограничена, потому что по неравенству Коши--Буняковского
	\[
	|(x_m,x_n)|
	\leqslant
	\|x_m\|\cdot\|x_n\|
	\leqslant
	R^2.
	\]
	Значит, из неё можно выделить сходящуюся подпоследовательность.
	
	Теперь применим диагональный метод Кантора.
	
	\begin{itemize}
		\item Из \(\{x_n\}\) выделим подпоследовательность, на которой сходится
		\[
		(x_1,x_n).
		\]
		
		\item Из неё выделим подпоследовательность, на которой сходится
		\[
		(x_2,x_n).
		\]
		
		\item Продолжая процесс, получаем систему вложенных подпоследовательностей.
		
		\item Берём диагональную подпоследовательность и обозначаем её
		\[
		y_k=x_{n_k}.
		\]
	\end{itemize}
	
	Тогда для каждого фиксированного \(m\in\mathbb{N}\) существует предел
	\[
	\lim_{k\to\infty}(x_m,y_k).
	\]
	
	\textbf{3. Построение функционала на плотном подпространстве.}
	
	Покажем, что для любого
	\[
	z\in L_0=\operatorname{span}\{x_1,x_2,\ldots\}
	\]
	существует предел
	\[
	\lim_{k\to\infty}(z,y_k).
	\]
	Действительно, если
	\[
	z=\sum_{m=1}^{N}\alpha_m x_m,
	\]
	то
	\[
	(z,y_k)
	=
	\sum_{m=1}^{N}\alpha_m(x_m,y_k)
	\]
	с точностью до выбранного соглашения о линейности скалярного произведения. Каждое
	\((x_m,y_k)\) имеет предел при \(k\to\infty\), значит, предел имеет и конечная линейная
	комбинация.
	
	Зададим на \(L_0\) функционал
	\[
	F(z)=\lim_{k\to\infty}(z,y_k).
	\]
	Он линеен, поскольку предел сохраняет линейные операции.
	
	Кроме того, \(F\) ограничен. Для любого \(z\in L_0\)
	\[
	|F(z)|
	=
	\left|\lim_{k\to\infty}(z,y_k)\right|
	\leqslant
	\limsup_{k\to\infty}|(z,y_k)|.
	\]
	По неравенству Коши--Буняковского
	\[
	|(z,y_k)|
	\leqslant
	\|z\|\cdot\|y_k\|
	\leqslant
	R\|z\|.
	\]
	Следовательно,
	\[
	|F(z)|\leqslant R\|z\|.
	\]
	Значит, \(F\) ограничен на \(L_0\).
	
	Так как
	\[
	L=\overline{L_0},
	\]
	функционал \(F\) единственным образом продолжается до ограниченного линейного функционала
	на всём \(L\). Это продолжение снова обозначим через \(F\). При этом
	\[
	F\in L^*,
	\qquad
	\|F\|\leqslant R.
	\]
	
	\textbf{4. Предъявление слабого предела по теореме Рисса--Фреше.}
	
	Так как \(L\) --- гильбертово пространство, по теореме Рисса--Фреше существует элемент
	\(x\in L\), такой что
	\[
	F(z)=(z,x)
	\qquad
	\forall z\in L.
	\]
	Причём
	\[
	\|x\|=\|F\|\leqslant R.
	\]
	Следовательно,
	\[
	x\in\overline{B}_R(0).
	\]
	
	Теперь докажем, что
	\[
	y_k\weakconv x
	\qquad
	\text{в } L.
	\]
	По построению \(F\)
	\[
	F(z)=\lim_{k\to\infty}(z,y_k),
	\]
	а по выбору \(x\)
	\[
	F(z)=(z,x).
	\]
	Значит,
	\[
	(z,y_k)\to(z,x)
	\qquad
	\forall z\in L.
	\]
	По сопряжённой симметрии скалярного произведения это равносильно
	\[
	(y_k,z)\to(x,z)
	\qquad
	\forall z\in L.
	\]
	Следовательно,
	\[
	y_k\weakconv x
	\qquad
	\text{в } L.
	\]
	
	\textbf{5. Переход от \(L\) ко всему \(H\).}
	
	Так как \(L\) --- замкнутое подпространство гильбертова пространства, по теореме о проекции
	\[
	H=L\oplus L^\perp.
	\]
	Поэтому любой \(h\in H\) единственным образом представляется в виде
	\[
	h=l+l^\perp,
	\qquad
	l\in L,
	\qquad
	l^\perp\in L^\perp.
	\]
	Так как
	\[
	y_k\in L,
	\qquad
	x\in L,
	\]
	то
	\[
	(y_k,l^\perp)=0,
	\qquad
	(x,l^\perp)=0.
	\]
	Следовательно,
	\[
	(y_k,h)
	=
	(y_k,l+l^\perp)
	=
	(y_k,l)
	\to
	(x,l)
	=
	(x,l+l^\perp)
	=
	(x,h).
	\]
	То есть
	\[
	(y_k,h)\to(x,h)
	\qquad
	\forall h\in H.
	\]
	По критерию слабой сходимости в гильбертовом пространстве
	\[
	y_k\weakconv x
	\qquad
	\text{в } H.
	\]
	
	При этом уже доказано, что
	\[
	x\in\overline{B}_R(0).
	\]
	Значит, из любой последовательности в шаре \(\overline{B}_R(0)\) можно выделить
	подпоследовательность, слабо сходящуюся к элементу того же шара. Поэтому замкнутый шар
	в гильбертовом пространстве слабо секвенциально компактен.
\end{proof}
\newpage
	
	% VIII. ОБРАТНЫЙ ОПЕРАТОР. ОБРАТИМОСТЬ
	\section{Обратный оператор. Обратимость}
	\hypertarget{bilet4}{}\setupnav{3}{5} \subsection{Обратимость линейного, ограниченного снизу \texorpdfstring{\((\|Ax\|\geqslant k\|x\|)\)}{(||Ax|| >= k||x||)} оператора}

\begin{definition}[Обратимость оператора на образе]
	Пусть \(E_1,E_2\) --- нормированные пространства,
	\[
	A\in\mathcal{L}(E_1,E_2).
	\]
	Оператор \(A\) называется \textbf{обратимым на своём образе}, если для любого
	\[
	y\in \operatorname{Im}A
	\]
	уравнение
	\[
	Ax=y
	\]
	имеет единственное решение \(x\in E_1\).
	
	В этом случае определён обратный оператор
	\[
	A^{-1}:\operatorname{Im}A\to E_1,
	\]
	задаваемый правилом
	\[
	A^{-1}y=x
	\quad\Longleftrightarrow\quad
	Ax=y.
	\]
\end{definition}

\begin{remark}[Левый и правый обратные]
	Для оператора
	\[
	A:E_1\to E_2
	\]
	можно различать левый и правый обратные операторы.
	
	\begin{itemize}
		\item \textbf{Левый обратный} отвечает за единственность решения:
		\[
		LA=I_{E_1}.
		\]
		
		\item \textbf{Правый обратный} отвечает за существование решения на образе:
		\[
		AR=I_{\operatorname{Im}A}.
		\]
	\end{itemize}
	
	Если оба оператора существуют и совпадают, то получаем обычный обратный оператор
	\[
	A^{-1}.
	\]
	В этом билете обратный оператор рассматривается именно как оператор
	\[
	A^{-1}:\operatorname{Im}A\to E_1.
	\]
\end{remark}

\begin{definition}[Ограниченный снизу оператор]
	Пусть \(E_1,E_2\) --- нормированные пространства,
	\[
	A\in\mathcal{L}(E_1,E_2).
	\]
	Оператор \(A\) называется \textbf{ограниченным снизу}, если существует число
	\[
	k>0,
	\]
	такое что
	\[
	\|Ax\|_{E_2}\geqslant k\|x\|_{E_1}
	\qquad
	\forall x\in E_1.
	\]
\end{definition}

\begin{lemma}[Ограниченность снизу влечёт инъективность]
	Если оператор
	\[
	A\in\mathcal{L}(E_1,E_2)
	\]
	ограничен снизу, то
	\[
	\ker A=\{0\}.
	\]
	Следовательно, \(A\) инъективен.
\end{lemma}

\begin{proof}
	Пусть
	\[
	x\in\ker A.
	\]
	Тогда \(Ax=0\). Так как \(A\) ограничен снизу, существует \(k>0\), такое что
	\[
	\|Ax\|\geqslant k\|x\|.
	\]
	Но \(Ax=0\), поэтому
	\[
	0=\|Ax\|\geqslant k\|x\|.
	\]
	Так как \(k>0\), получаем
	\[
	\|x\|=0.
	\]
	Значит,
	\[
	x=0.
	\]
	Следовательно,
	\[
	\ker A=\{0\}.
	\]
\end{proof}

\begin{theorem}[Критерий ограниченности обратного оператора]
	Пусть \(E_1,E_2\) --- нормированные пространства,
	\[
	A\in\mathcal{L}(E_1,E_2).
	\]
	Тогда обратный оператор
	\[
	A^{-1}:\operatorname{Im}A\to E_1
	\]
	существует и ограничен тогда и только тогда, когда оператор \(A\) ограничен снизу, то есть
	\[
	\exists k>0
	\quad
	\forall x\in E_1
	\qquad
	\|Ax\|\geqslant k\|x\|.
	\]
\end{theorem}

\begin{proof}
	Докажем обе импликации.
	
	\textbf{1. Если \(A^{-1}\) существует и ограничен, то \(A\) ограничен снизу.}
	
	Пусть
	\[
	A^{-1}\in\mathcal{L}(\operatorname{Im}A,E_1).
	\]
	Возьмём произвольный элемент \(x\in E_1\). Тогда \(Ax\in\operatorname{Im}A\), и по определению
	обратного оператора
	\[
	A^{-1}Ax=x.
	\]
	Следовательно,
	\[
	\|x\|
	=
	\|A^{-1}Ax\|
	\leqslant
	\|A^{-1}\|\cdot\|Ax\|.
	\]
	Отсюда
	\[
	\|Ax\|
	\geqslant
	\frac{1}{\|A^{-1}\|}\|x\|.
	\]
	Значит, \(A\) ограничен снизу с константой
	\[
	k=\frac{1}{\|A^{-1}\|}.
	\]
	
	\textbf{2. Если \(A\) ограничен снизу, то \(A^{-1}\) существует и ограничен.}
	
	Пусть существует \(k>0\), такое что
	\[
	\|Ax\|\geqslant k\|x\|
	\qquad
	\forall x\in E_1.
	\]
	Докажем, что
	\[
	A^{-1}:\operatorname{Im}A\to E_1
	\]
	корректно определён и ограничен.
	
	\begin{itemize}
		\item \textbf{Существование обратного оператора на образе.}
		
		По лемме из ограниченности снизу следует
		\[
		\ker A=\{0\}.
		\]
		Значит, \(A\) инъективен.
		
		Поэтому для каждого \(y\in\operatorname{Im}A\) существует единственный \(x\in E_1\), такой что
		\[
		Ax=y.
		\]
		Следовательно, корректно определён оператор
		\[
		A^{-1}:\operatorname{Im}A\to E_1.
		\]
		
		\item \textbf{Линейность обратного оператора.}
		
		Пусть
		\[
		y_1,y_2\in\operatorname{Im}A,
		\qquad
		\alpha,\beta\in\K.
		\]
		Выберем \(x_1,x_2\in E_1\), такие что
		\[
		Ax_1=y_1,
		\qquad
		Ax_2=y_2.
		\]
		Тогда
		\[
		A(\alpha x_1+\beta x_2)
		=
		\alpha Ax_1+\beta Ax_2
		=
		\alpha y_1+\beta y_2.
		\]
		По единственности прообраза
		\[
		A^{-1}(\alpha y_1+\beta y_2)
		=
		\alpha x_1+\beta x_2
		=
		\alpha A^{-1}y_1+\beta A^{-1}y_2.
		\]
		Значит, \(A^{-1}\) линеен.
		
		\item \textbf{Ограниченность обратного оператора.}
		
		Возьмём произвольный
		\[
		y\in\operatorname{Im}A.
		\]
		Тогда \(y=Ax\) для некоторого \(x\in E_1\), причём
		\[
		A^{-1}y=x.
		\]
		Используем ограниченность \(A\) снизу:
		\[
		\|y\|
		=
		\|Ax\|
		\geqslant
		k\|x\|
		=
		k\|A^{-1}y\|.
		\]
		Отсюда
		\[
		\|A^{-1}y\|
		\leqslant
		\frac{1}{k}\|y\|.
		\]
	\end{itemize}
	
	Итак,
	\[
	A^{-1}\in\mathcal{L}(\operatorname{Im}A,E_1),
	\qquad
	\|A^{-1}\|\leqslant \frac{1}{k}.
	\]
\end{proof}

\begin{remark}[Полезное следствие: замкнутость образа]
	Если \(E_1\) --- банахово пространство и оператор
	\[
	A\in\mathcal{L}(E_1,E_2)
	\]
	ограничен снизу, то
	\[
	\operatorname{Im}A
	\]
	замкнут в \(E_2\).
	
	Действительно, пусть
	\[
	y_n\in\operatorname{Im}A,
	\qquad
	y_n\to y.
	\]
	Представим
	\[
	y_n=Az_n.
	\]
	Так как \(\{y_n\}\) сходится, она фундаментальна. По оценке снизу
	\[
	\|y_n-y_m\|
	=
	\|Az_n-Az_m\|
	=
	\|A(z_n-z_m)\|
	\geqslant
	k\|z_n-z_m\|.
	\]
	Значит, \(\{z_n\}\) фундаментальна. Так как \(E_1\) банахово, существует \(z\in E_1\), такое что
	\[
	z_n\to z.
	\]
	По непрерывности \(A\)
	\[
	Az_n\to Az.
	\]
	Но \(Az_n=y_n\to y\), поэтому
	\[
	y=Az\in\operatorname{Im}A.
	\]
	Следовательно, \(\operatorname{Im}A\) замкнут.
\end{remark}
\newpage
	\hypertarget{bilet5}{}\setupnav{4}{6} \subsection{Обратимость возмущенного оператора \(A+\Delta A\)}

\begin{theorem}[Обратимость возмущенного оператора]
	\textbf{Условия:}
	\begin{itemize}
		\item $E$ --- банахово пространство;
		\item $A \in \mathcal{L}(E)$ --- обратимый оператор, причем
		$A^{-1} \in \mathcal{L}(E)$;
		\item $\Delta A \in \mathcal{L}(E)$;
		\item выполнено условие
		\[
		\|A^{-1}\|\,\|\Delta A\| < 1.
		\]
	\end{itemize}
	
	\textbf{Следовательно:}
	\begin{itemize}
		\item оператор $A+\Delta A$ обратим;
		\item $(A+\Delta A)^{-1} \in \mathcal{L}(E)$.
	\end{itemize}
\end{theorem}

\begin{proof}
	\begin{enumerate}
		\item \textbf{Сведение к оператору вида $I+B$.}
		Положим
		\[
		B := A^{-1}\Delta A \in \mathcal{L}(E).
		\]
		Тогда
		\[
		\|B\|
		\leqslant
		\|A^{-1}\|\,\|\Delta A\|
		<1,
		\]
		и
		\[
		A+\Delta A
		=
		A(I+A^{-1}\Delta A)
		=
		A(I+B).
		\]
		
		\item \textbf{Ряд Неймана.}
		Рассмотрим операторный ряд
		\[
		S := \sum_{n=0}^{\infty}(-B)^n.
		\]
		Так как
		\[
		\|B^n\|\leqslant\|B\|^n
		\]
		и числовой ряд
		\[
		\sum_{n=0}^{\infty}\|B\|^n
		\]
		сходится, то ряд $\sum\limits_{n=0}^{\infty}(-B)^n$
		сходится в банаховом пространстве $\mathcal{L}(E)$.
		Следовательно, $S\in\mathcal{L}(E)$.
		
		\item \textbf{Обратимость $I+B$.}
		Для частичной суммы
		\[
		S_N=\sum_{n=0}^{N}(-B)^n
		\]
		имеем
		\[
		(I+B)S_N=S_N(I+B)=I+(-1)^N B^{N+1}.
		\]
		Поскольку
		\[
		\|B^{N+1}\|\leqslant\|B\|^{N+1}\to0,
		\]
		при $N\to\infty$ получаем
		\[
		(I+B)S=S(I+B)=I.
		\]
		Значит,
		\[
		(I+B)^{-1}=S\in\mathcal{L}(E).
		\]
		
		\item \textbf{Обратимость $A+\Delta A$.}
		Так как $A$ и $I+B$ обратимы,
		\[
		(A+\Delta A)^{-1}
		=
		(I+B)^{-1}A^{-1}.
		\]
		Причем оба множителя принадлежат $\mathcal{L}(E)$, поэтому
		\[
		(A+\Delta A)^{-1}\in\mathcal{L}(E).
		\]
	\end{enumerate}
\end{proof}

\begin{corollary}
	В условиях теоремы
	\[
	(A+\Delta A)^{-1}
	=
	\sum_{n=0}^{\infty}
	(-A^{-1}\Delta A)^n A^{-1}.
	\]
	Кроме того,
	\[
	\|(A+\Delta A)^{-1}\|
	\leqslant
	\frac{\|A^{-1}\|}
	{1-\|A^{-1}\|\,\|\Delta A\|}.
	\]
\end{corollary}

\begin{proof}
	Из доказательства теоремы
	\[
	(I+B)^{-1}
	=
	\sum_{n=0}^{\infty}(-B)^n,
	\qquad
	B=A^{-1}\Delta A.
	\]
	Поэтому
	\[
	(A+\Delta A)^{-1}
	=
	(I+B)^{-1}A^{-1}
	=
	\sum_{n=0}^{\infty}
	(-A^{-1}\Delta A)^n A^{-1}.
	\]
	
	Также
	\[
	\|(I+B)^{-1}\|
	\leqslant
	\sum_{n=0}^{\infty}\|B\|^n
	=
	\frac{1}{1-\|B\|}.
	\]
	Следовательно,
	\[
	\|(A+\Delta A)^{-1}\|
	\leqslant
	\frac{\|A^{-1}\|}
	{1-\|A^{-1}\|\,\|\Delta A\|}.
	\]
\end{proof}

\begin{idea}
	\textbf{Суть:} Обратимость устойчива относительно достаточно малых
	возмущений. Если $A$ обратим и
	\[
	\|\Delta A\| < \frac{1}{\|A^{-1}\|},
	\]
	то $A+\Delta A$ также имеет ограниченный обратный оператор.
\end{idea}\newpage
	\hypertarget{bilet6}{}\setupnav{5}{7} \providecommand{\K}{\mathbb{K}}

\subsection{Формулировка теоремы Банаха об обратном операторе.
	Доказательство теоремы Банаха в случае гильбертова пространства}

\begin{theorem}[Банаха об обратном операторе]
	Пусть \(E_1,E_2\) --- банаховы пространства,
	\[
	A\in\mathcal{L}(E_1,E_2),
	\]
	и \(A:E_1\to E_2\) --- биекция. Тогда
	\[
	A^{-1}\in\mathcal{L}(E_2,E_1).
	\]
\end{theorem}

\begin{theorem}[Гильбертов случай]
	Пусть \(H\) --- гильбертово пространство,
	\[
	A\in\mathcal{L}(H),
	\]
	и \(A:H\to H\) --- биекция. Тогда
	\[
	A^{-1}\in\mathcal{L}(H).
	\]
\end{theorem}

\begin{proof}
	Сначала получим ограниченный обратный к \(A^*\), а затем применим то же рассуждение
	к оператору \(A^*\).
	
	\textbf{1. Докажем, что \(A^*\) ограничен снизу.}
	
	По билету 8.1 для этого достаточно получить оценку
	\[
	\exists m>0
	\quad
	\forall x\in H
	\qquad
	\|A^*x\|\geqslant m\|x\|.
	\]
	
	Оценка однородна по \(x\), поэтому достаточно рассмотреть векторы, для которых
	\[
	\|A^*x\|=1.
	\]
	Положим
	\[
	S=\{x\in H:\|A^*x\|=1\}.
	\]
	Тогда достаточно доказать, что \(S\) ограничено:
	\[
	\exists C>0
	\quad
	\forall x\in S
	\qquad
	\|x\|\leqslant C.
	\]
	
	По теореме Хана достаточно показать, что \(S\) слабо ограничено.
	
	Пусть \(y\in H\). Так как \(A\) сюръективен, существует \(z\in H\), такое что
	\[
	Az=y.
	\]
	Тогда для любого \(x\in S\)
	\[
	|(x,y)|
	=
	|(x,Az)|
	=
	|(A^*x,z)|
	\leqslant
	\|A^*x\|\|z\|
	=
	\|z\|.
	\]
	Значит, \(S\) слабо ограничено, а по теореме Хана --- ограничено:
	\[
	\exists C>0
	\quad
	\forall x\in S
	\qquad
	\|x\|\leqslant C.
	\]
	
	Теперь перейдём от ограниченности \(S\) к оценке снизу.
	Для \(u\neq0\) сначала проверим, что \(A^*u\neq0\).
	
	Так как \(A\) сюръективен,
	\[
	\operatorname{Im}A=H,
	\]
	поэтому
	\[
	\ker A^*
	=
	(\operatorname{Im}A)^\perp
	=
	\{0\}.
	\]
	Следовательно,
	\[
	\frac{u}{\|A^*u\|}\in S.
	\]
	Отсюда
	\[
	\frac{\|u\|}{\|A^*u\|}
	\leqslant C,
	\]
	то есть
	\[
	\|A^*u\|
	\geqslant
	\frac1C\|u\|.
	\]
	Значит, \(A^*\) ограничен снизу.
	
	\textbf{2. Покажем, что \(A^*\) --- биекция.}
	
	Инъективность уже получена:
	\[
	\ker A^*=\{0\}.
	\]
	
	Осталось доказать сюръективность. Так как \(A\) инъективен,
	\[
	\ker A=\{0\}.
	\]
	По разложению
	\[
	H=\overline{\operatorname{Im}A^*}\oplus\ker A
	\]
	получаем
	\[
	\overline{\operatorname{Im}A^*}=H.
	\]
	
	С другой стороны, \(A^*\) ограничен снизу, поэтому по билету 8.1
	\[
	\operatorname{Im}A^*
	\]
	замкнут. Следовательно,
	\[
	\operatorname{Im}A^*
	=
	\overline{\operatorname{Im}A^*}
	=
	H.
	\]
	
	Таким образом, \(A^*\) --- биекция. Так как \(A^*\) ограничен снизу,
	по билету 8.1
	\[
	(A^*)^{-1}\in\mathcal L(H).
	\]
	
	\textbf{3. Применим доказанное к \(A^*\).}
	
	Оператор
	\[
	A^*\in\mathcal L(H)
	\]
	является биекцией. Поэтому, повторяя доказанное выше для \(A^*\), получаем
	\[
	((A^*)^*)^{-1}\in\mathcal L(H).
	\]
	Так как
	\[
	(A^*)^*=A,
	\]
	имеем
	\[
	\boxed{A^{-1}\in\mathcal L(H)}.
	\]
\end{proof}
\newpage
	
	% IX. СПЕКТР. РЕЗОЛЬВЕНТА
	\section{Спектр. Резольвента}
	\hypertarget{bilet7}{}\setupnav{6}{8} \subsection{Операторнозначные функции комплексного переменного. Аналитичность резольвенты. Спектральный радиус. Основная теорема о спектре}

\begin{definition}[Операторнозначная функция комплексного переменного]
	Пусть \(E\) --- банахово пространство над \(\mathbb{C}\), а \(D\subset\mathbb{C}\).
	Отображение
	\[
	F:D\to\mathcal{L}(E)
	\]
	называется \textbf{операторнозначной функцией комплексного переменного}.
\end{definition}

\begin{definition}[Аналитичность операторнозначной функции]
	Функция
	\[
	F:D\to\mathcal{L}(E)
	\]
	называется \textbf{аналитической} в точке \(\lambda_0\in D\), если существует предел
	\[
	F'(\lambda_0)
	=
	\lim_{\lambda\to\lambda_0}
	\frac{F(\lambda)-F(\lambda_0)}{\lambda-\lambda_0}
	\]
	в норме пространства \(\mathcal{L}(E)\).
	
	Если \(F\) аналитична в каждой точке области \(D\), то \(F\) называется аналитической в \(D\).
\end{definition}


\begin{definition}[Резольвентное множество, резольвента и спектр]
	Пусть \(E\) --- комплексное банахово пространство,
	\[
	A\in\mathcal{L}(E).
	\]
	Для \(\lambda\in\mathbb{C}\) обозначим
	\[
	A_\lambda=A-\lambda I.
	\]
	
	Число \(\lambda\in\mathbb{C}\) называется \textbf{регулярным значением} оператора \(A\), если оператор
	\[
	A_\lambda=A-\lambda I
	\]
	непрерывно обратим, то есть
	\[
	A_\lambda^{-1}\in\mathcal{L}(E).
	\]
	
	\textbf{Резольвентным множеством} называется множество
	\[
	\rho(A)=\{\lambda\in\mathbb{C}: A_\lambda^{-1}\in\mathcal{L}(E)\}.
	\]
	
	Для \(\lambda\in\rho(A)\) оператор
	\[
	R_\lambda=R_\lambda(A)=(A-\lambda I)^{-1}
	\]
	называется \textbf{резольвентой} оператора \(A\).
	
	\textbf{Спектром} оператора \(A\) называется множество
	\[
	\sigma(A)=\mathbb{C}\setminus\rho(A).
	\]
\end{definition}

\begin{definition}[Спектральный радиус]
	Спектральным радиусом оператора \(A\) называется число
	\[
	r(A)=\sup_{\lambda\in\sigma(A)}|\lambda|.
	\]
\end{definition}

\begin{proposition}[Резольвента вне круга радиуса \(\|A\|\)]
	Если
	\[
	|\lambda|>\|A\|,
	\]
	то
	\[
	\lambda\in\rho(A),
	\]
	причём
	\[
	R_\lambda
	=
	-\frac{1}{\lambda}
	\sum_{n=0}^{\infty}\frac{A^n}{\lambda^n}.
	\]
\end{proposition}

\begin{proof}
	\textbf{1. Сведение к ряду Неймана.}
	
	Имеем
	\[
	A-\lambda I
	=
	-\lambda\left(I-\frac{1}{\lambda}A\right).
	\]
	Так как
	\[
	\left\|\frac{1}{\lambda}A\right\|
	=
	\frac{\|A\|}{|\lambda|}
	<1,
	\]
	то по теореме о ряде Неймана оператор
	\[
	I-\frac{1}{\lambda}A
	\]
	обратим, и
	\[
	\left(I-\frac{1}{\lambda}A\right)^{-1}
	=
	\sum_{n=0}^{\infty}\frac{A^n}{\lambda^n}.
	\]
	
	\textbf{2. Формула для резольвенты.}
	
	Из разложения
	\[
	A-\lambda I
	=
	-\lambda\left(I-\frac{1}{\lambda}A\right)
	\]
	получаем
	\[
	(A-\lambda I)^{-1}
	=
	-\frac{1}{\lambda}
	\left(I-\frac{1}{\lambda}A\right)^{-1}.
	\]
	Следовательно,
	\[
	R_\lambda
	=
	-\frac{1}{\lambda}
	\sum_{n=0}^{\infty}\frac{A^n}{\lambda^n}.
	\]
	Значит,
	\[
	\lambda\in\rho(A).
	\]
\end{proof}

\begin{corollary}
	Спектр оператора \(A\) лежит в замкнутом круге:
	\[
	\sigma(A)\subset\{\lambda\in\mathbb{C}:|\lambda|\leqslant \|A\|\}.
	\]
\end{corollary}

\begin{theorem}[Аналитичность резольвенты]
	Резольвентное множество \(\rho(A)\) открыто, а функция
	\[
	R_\lambda=(A-\lambda I)^{-1}
	\]
	аналитична на \(\rho(A)\). Более того,
	\[
	R'_\lambda=R_\lambda^2.
	\]
\end{theorem}

\begin{proof}
	\textbf{1. Открытость резольвентного множества.}
	
	Возьмём произвольную точку
	\[
	\lambda_0\in\rho(A).
	\]
	Тогда оператор
	\[
	R_{\lambda_0}=(A-\lambda_0I)^{-1}
	\]
	существует и ограничен.
	
	Пусть \(\lambda\) близко к \(\lambda_0\). Обозначим
	\[
	\delta=\lambda-\lambda_0.
	\]
	Тогда
	\[
	A-\lambda I
	=
	A-\lambda_0I-\delta I
	=
	(A-\lambda_0I)(I-\delta R_{\lambda_0}).
	\]
	Если
	\[
	|\delta|\mathbb{C}dot\|R_{\lambda_0}\|<1,
	\]
	то по теореме о ряде Неймана оператор
	\[
	I-\delta R_{\lambda_0}
	\]
	обратим. Значит, обратим и оператор \(A-\lambda I\).
	
	Следовательно, при
	\[
	|\lambda-\lambda_0|<\frac{1}{\|R_{\lambda_0}\|}
	\]
	имеем
	\[
	\lambda\in\rho(A).
	\]
	Значит, \(\rho(A)\) открыто.
	
	\textbf{2. Локальное разложение резольвенты.}
	
	Из равенства
	\[
	A-\lambda I
	=
	(A-\lambda_0I)(I-(\lambda-\lambda_0)R_{\lambda_0})
	\]
	получаем
	\[
	R_\lambda
	=
	(I-(\lambda-\lambda_0)R_{\lambda_0})^{-1}R_{\lambda_0}.
	\]
	При
	\[
	|\lambda-\lambda_0|\mathbb{C}dot\|R_{\lambda_0}\|<1
	\]
	по ряду Неймана:
	\[
	(I-(\lambda-\lambda_0)R_{\lambda_0})^{-1}
	=
	\sum_{n=0}^{\infty}(\lambda-\lambda_0)^nR_{\lambda_0}^n.
	\]
	Следовательно,
	\[
	R_\lambda
	=
	\sum_{n=0}^{\infty}(\lambda-\lambda_0)^nR_{\lambda_0}^{n+1}.
	\]
	Это степенной ряд по \(\lambda-\lambda_0\) со значениями в \(\mathcal{L}(E)\). Поэтому
	\[
	\lambda\mapsto R_\lambda
	\]
	аналитична в окрестности \(\lambda_0\).
	
	\textbf{3. Формула для производной.}
	
	Из локального разложения сразу получаем
	\[
	R'_{\lambda_0}=R_{\lambda_0}^2.
	\]
	Так как \(\lambda_0\in\rho(A)\) была произвольной,
	\[
	R'_\lambda=R_\lambda^2
	\qquad
	\forall\lambda\in\rho(A).
	\]
\end{proof}

\begin{remark}[Равенство Гильберта для резольвенты]
	Если
	\[
	\lambda,\mu\in\rho(A),
	\]
	то
	\[
	R_\lambda-R_\mu
	=
	(\lambda-\mu)R_\lambda R_\mu.
	\]
	
	Действительно,
	\[
	R_\lambda-R_\mu
	=
	R_\lambda R_\mu^{-1}R_\mu
	-
	R_\lambda R_\lambda^{-1}R_\mu
	=
	R_\lambda(R_\mu^{-1}-R_\lambda^{-1})R_\mu.
	\]
	Но
	\[
	R_\mu^{-1}=A-\mu I,
	\qquad
	R_\lambda^{-1}=A-\lambda I.
	\]
	Значит,
	\[
	R_\mu^{-1}-R_\lambda^{-1}
	=
	(A-\mu I)-(A-\lambda I)
	=
	(\lambda-\mu)I.
	\]
	Следовательно,
	\[
	R_\lambda-R_\mu
	=
	(\lambda-\mu)R_\lambda R_\mu.
	\]
\end{remark}

\begin{theorem}[Основная теорема о спектре]
	Пусть \(E\) --- комплексное банахово пространство и
	\[
	A\in\mathcal{L}(E).
	\]
	Тогда:
	\begin{itemize}
		\item спектр \(\sigma(A)\) --- непустой компакт;
		\item спектральный радиус задаётся формулой
		\[
		r(A)
		=
		\lim_{n\to\infty}\sqrt[n]{\|A^n\|}.
		\]
	\end{itemize}
\end{theorem}

\begin{proof}
	\textbf{1. Замкнутость и ограниченность спектра.}
	
	Из предыдущего утверждения знаем, что
	\[
	\rho(A)
	\]
	открыто. Поэтому его дополнение
	\[
	\sigma(A)=\mathbb{C}\setminus\rho(A)
	\]
	замкнуто.
	
	Кроме того, если
	\[
	|\lambda|>\|A\|,
	\]
	то
	\[
	\lambda\in\rho(A).
	\]
	Значит,
	\[
	\sigma(A)\subset\{\lambda\in\mathbb{C}:|\lambda|\leqslant\|A\|\}.
	\]
	Следовательно, \(\sigma(A)\) замкнуто и ограничено. Если оно непусто, то оно компактно.
	
	Осталось доказать непустоту.
	
	\textbf{2. Непустота спектра.}
	
	Докажем от противного. Пусть
	\[
	\sigma(A)=\varnothing.
	\]
	Тогда
	\[
	\rho(A)=\mathbb{C},
	\]
	то есть резольвента
	\[
	R_\lambda=(A-\lambda I)^{-1}
	\]
	определена и аналитична на всей комплексной плоскости.
	
	При \(|\lambda|>\|A\|\) уже есть формула
	\[
	R_\lambda
	=
	-\frac{1}{\lambda}
	\sum_{n=0}^{\infty}\frac{A^n}{\lambda^n}.
	\]
	Поэтому
	\[
	\|R_\lambda\|
	\leqslant
	\frac{1}{|\lambda|}
	\sum_{n=0}^{\infty}
	\frac{\|A\|^n}{|\lambda|^n}
	=
	\frac{1}{|\lambda|}
	\mathbb{C}dot
	\frac{1}{1-\frac{\|A\|}{|\lambda|}}
	=
	\frac{1}{|\lambda|-\|A\|}.
	\]
	Отсюда
	\[
	\|R_\lambda\|\to0
	\qquad
	\text{при }|\lambda|\to\infty.
	\]
	
	Значит, \(R_\lambda\) ограничена вне большого круга. Внутри этого круга она ограничена по
	непрерывности на компакте. Следовательно, \(R_\lambda\) ограничена на всей \(\mathbb{C}\).
	
	По теореме Лиувилля для операторнозначных функций \(R_\lambda\) постоянна. Так как
	\[
	R_\lambda\to0
	\qquad
	|\lambda|\to\infty,
	\]
	эта постоянная равна нулю:
	\[
	R_\lambda\equiv0.
	\]
	Но \(R_\lambda\) является обратным оператором к \(A-\lambda I\), значит не может быть нулевым
	оператором. Получили противоречие.
	
	Следовательно,
	\[
	\sigma(A)\ne\varnothing.
	\]
	Значит, \(\sigma(A)\) --- непустой компакт.
	
	\textbf{3. Формула спектрального радиуса.}
	
	Обозначим
	\[
	L=\lim_{n\to\infty}\sqrt[n]{\|A^n\|}.
	\]
	Предел существует по стандартной лемме для субмультипликативной последовательности,
	поскольку
	\[
	\|A^{n+m}\|\leqslant\|A^n\|\mathbb{C}dot\|A^m\|.
	\]
	
	Докажем, что
	\[
	r(A)=L.
	\]
	
	\begin{itemize}
		\item \textbf{Оценка \(r(A)\leqslant L\).}
		
		Пусть
		\[
		\lambda\in\sigma(A).
		\]
		Тогда для любого \(n\in\mathbb{N}\)
		\[
		\lambda^n\in\sigma(A^n).
		\]
		Действительно, если бы \(A^n-\lambda^nI\) был обратим, то из разложения
		\[
		A^n-\lambda^nI
		=
		(A-\lambda I)
		(A^{n-1}+\lambda A^{n-2}+\ldots+\lambda^{n-1}I)
		\]
		и коммутативности множителей следовала бы обратимость \(A-\lambda I\), что противоречит
		\[
		\lambda\in\sigma(A).
		\]
		
		Так как спектр любого ограниченного оператора лежит в круге радиуса его нормы, то
		\[
		|\lambda|^n
		=
		|\lambda^n|
		\leqslant
		\|A^n\|.
		\]
		Следовательно,
		\[
		|\lambda|
		\leqslant
		\sqrt[n]{\|A^n\|}
		\qquad
		\forall n.
		\]
		Переходя к пределу по \(n\), получаем
		\[
		|\lambda|\leqslant L.
		\]
		Так как \(\lambda\in\sigma(A)\) произвольно,
		\[
		r(A)\leqslant L.
		\]
		
		\item \textbf{Оценка \(L\leqslant r(A)\).}
		
		Для \(|\lambda|>\|A\|\) мы знаем разложение
		\[
		R_\lambda
		=
		-\frac{1}{\lambda}
		\sum_{n=0}^{\infty}
		\frac{A^n}{\lambda^n}.
		\]
		Это ряд Лорана резольвенты в окрестности бесконечности.
		
		Так как резольвента аналитична на \(\rho(A)\), а
		\[
		\sigma(A)\subset\{\lambda:|\lambda|\leqslant r(A)\},
		\]
		то \(R_\lambda\) аналитична во всей внешности круга
		\[
		|\lambda|>r(A).
		\]
		По единственности ряда Лорана указанное разложение является лорановским разложением
		резольвенты во всей этой внешней области.
		
		Значит, ряд
		\[
		\sum_{n=0}^{\infty}
		\frac{A^n}{\lambda^{n+1}}
		\]
		сходится при каждом
		\[
		|\lambda|>r(A).
		\]
		Возьмём любое число \(s>r(A)\) и положим \(|\lambda|=s\). Из сходимости ряда следует, что
		его общий член стремится к нулю:
		\[
		\left\|
		\frac{A^n}{\lambda^{n+1}}
		\right\|
		=
		\frac{\|A^n\|}{s^{n+1}}
		\to0.
		\]
		Отсюда
		\[
		\limsup_{n\to\infty}\sqrt[n]{\frac{\|A^n\|}{s^{n+1}}}
		\leqslant1.
		\]
		То есть
		\[
		L\leqslant s.
		\]
		Так как это верно для любого \(s>r(A)\), получаем
		\[
		L\leqslant r(A).
		\]
	\end{itemize}
	
	Из двух оценок следует
	\[
	r(A)=L.
	\]
	То есть
	\[
	r(A)
	=
	\lim_{n\to\infty}\sqrt[n]{\|A^n\|}.
	\]
\end{proof}

\begin{idea}[Скелет билета]
	\begin{itemize}
		\item Для \(\lambda\in\mathbb{C}\) рассматриваем
		\[
		A_\lambda=A-\lambda I.
		\]
		
		\item Резольвентное множество:
		\[
		\rho(A)=\{\lambda: A_\lambda^{-1}\in\mathcal{L}(E)\}.
		\]
		
		\item Резольвента:
		\[
		R_\lambda=(A-\lambda I)^{-1}.
		\]
		
		\item Спектр:
		\[
		\sigma(A)=\mathbb{C}\setminus\rho(A).
		\]
		
		\item Если \(|\lambda|>\|A\|\), то
		\[
		R_\lambda
		=
		-\frac{1}{\lambda}
		\sum_{n=0}^{\infty}\frac{A^n}{\lambda^n}.
		\]
		Значит,
		\[
		\sigma(A)\subset\{|\lambda|\leqslant\|A\|\}.
		\]
		
		\item \(\rho(A)\) открыто, потому что обратимость сохраняется при малом возмущении:
		\[
		A-\lambda I
		=
		(A-\lambda_0I)(I-(\lambda-\lambda_0)R_{\lambda_0}).
		\]
		
		\item Резольвента аналитична, потому что локально раскладывается в ряд Неймана:
		\[
		R_\lambda
		=
		\sum_{n=0}^{\infty}
		(\lambda-\lambda_0)^nR_{\lambda_0}^{n+1}.
		\]
		
		\item Основная теорема:
		\[
		\sigma(A)\text{ --- непустой компакт},
		\qquad
		r(A)=\lim_{n\to\infty}\sqrt[n]{\|A^n\|}.
		\]
	\end{itemize}
\end{idea}\newpage
	
	% X. СОПРЯЖЕННЫЙ ОПЕРАТОР
	\section{Сопряженный оператор}
	\hypertarget{bilet8}{}\setupnav{7}{9} \subsection{Норма сопряженного оператора (в линейном нормированном пространстве)}

\begin{definition}[Сопряженное пространство]
	Пусть \(E\) --- линейное нормированное пространство над полем \(\mathbb{K}\).
	Сопряженным пространством к \(E\) называется пространство всех линейных ограниченных
	функционалов на \(E\):
	\[
	E^*=\mathcal{L}(E,\mathbb{K}).
	\]
\end{definition}

\begin{definition}[Сопряженный оператор]
	Пусть \(E_1,E_2\) --- линейные нормированные пространства и
	\[
	A\in\mathcal{L}(E_1,E_2).
	\]
	Сопряженным к \(A\) называется оператор
	\[
	A^*:E_2^*\to E_1^*,
	\]
	который каждому функционалу \(g\in E_2^*\) сопоставляет функционал
	\[
	A^*g\in E_1^*
	\]
	по правилу
	\[
	(A^*g)(x)=g(Ax),
	\qquad
	x\in E_1.
	\]
	То есть
	\[
	A^*g=g\circ A.
	\]
\end{definition}

\begin{remark}[Следствие Хана--Банаха]
	Для любого элемента \(y\) линейного нормированного пространства \(E\) выполнено
	\[
	\|y\|
	=
	\sup_{\|f\|_{E^*}\leqslant 1}|f(y)|.
	\]
	Именно этот факт будет нужен для доказательства неравенства
	\[
	\|A\|\leqslant \|A^*\|.
	\]
\end{remark}

\begin{theorem}[Норма сопряженного оператора]
	Пусть \(E_1,E_2\) --- линейные нормированные пространства и
	\[
	A\in\mathcal{L}(E_1,E_2).
	\]
	Тогда сопряженный оператор
	\[
	A^*:E_2^*\to E_1^*
	\]
	существует, является линейным ограниченным оператором и
	\[
	\|A^*\|=\|A\|.
	\]
\end{theorem}

\begin{proof}
	Докажем утверждение по смысловым шагам.
	
	\textbf{1. Корректность определения \(A^*g\).}
	
	Пусть
	\[
	g\in E_2^*.
	\]
	Определим отображение
	\[
	A^*g:E_1\to\mathbb{K}
	\]
	по правилу
	\[
	(A^*g)(x)=g(Ax).
	\]
	Проверим, что
	\[
	A^*g\in E_1^*.
	\]
	
	\begin{itemize}
		\item \textbf{Линейность функционала \(A^*g\).}
		
		Пусть
		\[
		x_1,x_2\in E_1,
		\qquad
		\alpha,\beta\in\mathbb{K}.
		\]
		Тогда
		\[
		(A^*g)(\alpha x_1+\beta x_2)
		=
		g(A(\alpha x_1+\beta x_2)).
		\]
		Так как \(A\) линеен,
		\[
		A(\alpha x_1+\beta x_2)
		=
		\alpha Ax_1+\beta Ax_2.
		\]
		Так как \(g\) линеен,
		\[
		g(\alpha Ax_1+\beta Ax_2)
		=
		\alpha g(Ax_1)+\beta g(Ax_2).
		\]
		Следовательно,
		\[
		(A^*g)(\alpha x_1+\beta x_2)
		=
		\alpha(A^*g)(x_1)+\beta(A^*g)(x_2).
		\]
		Значит, \(A^*g\) линеен.
		
		\item \textbf{Ограниченность функционала \(A^*g\).}
		
		Для любого \(x\in E_1\) имеем
		\[
		|(A^*g)(x)|
		=
		|g(Ax)|
		\leqslant
		\|g\|\cdot\|Ax\|.
		\]
		Так как \(A\) ограничен,
		\[
		\|Ax\|\leqslant \|A\|\cdot\|x\|.
		\]
		Поэтому
		\[
		|(A^*g)(x)|
		\leqslant
		\|g\|\cdot\|A\|\cdot\|x\|.
		\]
		Значит,
		\[
		A^*g\in E_1^*.
		\]
	\end{itemize}
	
	\textbf{2. Линейность и ограниченность оператора \(A^*\).}
	
	Теперь рассматриваем \(A^*\) уже как оператор
	\[
	A^*:E_2^*\to E_1^*.
	\]
	
	\begin{itemize}
		\item \textbf{Линейность оператора \(A^*\).}
		
		Пусть
		\[
		g_1,g_2\in E_2^*,
		\qquad
		\alpha,\beta\in\mathbb{K}.
		\]
		Тогда для любого \(x\in E_1\)
		\[
		(A^*(\alpha g_1+\beta g_2))(x)
		=
		(\alpha g_1+\beta g_2)(Ax).
		\]
		Отсюда
		\[
		(A^*(\alpha g_1+\beta g_2))(x)
		=
		\alpha g_1(Ax)+\beta g_2(Ax)
		=
		\alpha(A^*g_1)(x)+\beta(A^*g_2)(x).
		\]
		Так как равенство верно для всех \(x\in E_1\), получаем
		\[
		A^*(\alpha g_1+\beta g_2)
		=
		\alpha A^*g_1+\beta A^*g_2.
		\]
		Значит, \(A^*\) линеен.
		
		\item \textbf{Ограниченность оператора \(A^*\).}
		
		Из оценки первого пункта
		\[
		|(A^*g)(x)|
		\leqslant
		\|g\|\cdot\|A\|\cdot\|x\|
		\]
		получаем
		\[
		\|A^*g\|
		\leqslant
		\|A\|\cdot\|g\|.
		\]
		Следовательно,
		\[
		A^*\in\mathcal{L}(E_2^*,E_1^*)
		\]
		и
		\[
		\|A^*\|\leqslant\|A\|.
		\]
	\end{itemize}
	
	\textbf{3. Обратное неравенство \(\|A\|\leqslant\|A^*\|\).}
	
	Возьмём произвольный элемент
	\[
	x\in E_1.
	\]
	По следствию Хана--Банаха для элемента \(Ax\in E_2\)
	\[
	\|Ax\|
	=
	\sup_{\|g\|_{E_2^*}\leqslant1}|g(Ax)|.
	\]
	По определению сопряженного оператора
	\[
	g(Ax)=(A^*g)(x).
	\]
	Поэтому
	\[
	\|Ax\|
	=
	\sup_{\|g\|_{E_2^*}\leqslant1}|(A^*g)(x)|.
	\]
	Для каждого \(g\in E_2^*\), такого что \(\|g\|\leqslant1\), имеем
	\[
	|(A^*g)(x)|
	\leqslant
	\|A^*g\|\cdot\|x\|.
	\]
	Но
	\[
	\|A^*g\|
	\leqslant
	\|A^*\|\cdot\|g\|
	\leqslant
	\|A^*\|.
	\]
	Значит,
	\[
	|(A^*g)(x)|
	\leqslant
	\|A^*\|\cdot\|x\|.
	\]
	Переходя к супремуму по всем \(\|g\|\leqslant1\), получаем
	\[
	\|Ax\|
	\leqslant
	\|A^*\|\cdot\|x\|.
	\]
	Так как это верно для любого \(x\in E_1\), то
	\[
	\|A\|\leqslant\|A^*\|.
	\]
	
	\textbf{4. Равенство норм.}
	
	Мы доказали две оценки:
	\[
	\|A^*\|\leqslant\|A\|
	\]
	и
	\[
	\|A\|\leqslant\|A^*\|.
	\]
	Следовательно,
	\[
	\|A^*\|=\|A\|.
	\]
\end{proof}
\newpage
	\hypertarget{bilet9}{}\setupnav{8}{10} \subsection{Сопряженные операторы в гильбертовом пространстве. Равенство \texorpdfstring{\(H=\overline{\operatorname{Im} A}\oplus\ker A^*\)}{H = cl Im A plus ker A*}}

\begin{definition}[Сопряженный оператор в гильбертовом пространстве]
	Пусть \(H_1,H_2\) --- гильбертовы пространства и
	\[
	A\in\mathcal{L}(H_1,H_2).
	\]
	Оператор
	\[
	A^*:H_2\to H_1
	\]
	называется \textbf{сопряженным} к \(A\), если
	\[
	(Ax,y)_{H_2}=(x,A^*y)_{H_1}
	\qquad
	\forall x\in H_1,\ \forall y\in H_2.
	\]
\end{definition}

\begin{remark}[Свойства сопряженного оператора]
	Сопряженный оператор в гильбертовом пространстве существует и единственен. Также
	выполнены свойства:
	\[
	I^*=I,
	\]
	\[
	(\alpha A+\beta B)^*
	=
	\overline{\alpha}A^*+\overline{\beta}B^*,
	\]
	\[
	(AB)^*=B^*A^*.
	\]
	Они доказываются прямой проверкой из определения.
\end{remark}

\begin{theorem}[Разложение гильбертова пространства]
	Пусть \(H\) --- гильбертово пространство над \(\mathbb{K}\), и
	\[
	A\in\mathcal{L}(H).
	\]
	Тогда
	\[
	H=\overline{\operatorname{Im}A}\oplus\ker A^*.
	\]
\end{theorem}

\begin{proof}
	Докажем разложение через ортогональное дополнение к образу.
	
	\textbf{1. Простая часть: докажем равенство \((\operatorname{Im}A)^\perp=\ker A^*\).}
	
	Пусть
	\[
	y\in(\operatorname{Im}A)^\perp.
	\]
	Это означает, что
	\[
	(Az,y)=0
	\qquad
	\forall z\in H.
	\]
	По определению сопряженного оператора
	\[
	(Az,y)=(z,A^*y).
	\]
	Значит,
	\[
	(z,A^*y)=0
	\qquad
	\forall z\in H.
	\]
	Отсюда
	\[
	A^*y=0.
	\]
	Следовательно,
	\[
	y\in\ker A^*.
	\]
	
	Обратно, если
	\[
	y\in\ker A^*,
	\]
	то
	\[
	A^*y=0.
	\]
	Тогда для любого \(z\in H\)
	\[
	(Az,y)=(z,A^*y)=(z,0)=0.
	\]
	Значит,
	\[
	y\in(\operatorname{Im}A)^\perp.
	\]
	
	Итак,
	\[
	(\operatorname{Im}A)^\perp=\ker A^*.
	\]
	
	\textbf{2. Сложная часть: применим теорему о проекции.}
	
	Знаем, что для любого множества \(S\subset H\)
	\[
	S^\perp=\left(\overline{S}\right)^\perp.
	\]
	Поэтому
	\[
	(\operatorname{Im}A)^\perp
	=
	\left(\overline{\operatorname{Im}A}\right)^\perp.
	\]
	Так как
	\[
	\overline{\operatorname{Im}A}
	\]
	--- замкнутое подпространство в \(H\), то по теореме Рисса о проекции
	\[
	H
	=
	\overline{\operatorname{Im}A}
	\oplus
	\left(\overline{\operatorname{Im}A}\right)^\perp.
	\]
	Но
	\[
	\left(\overline{\operatorname{Im}A}\right)^\perp
	=
	(\operatorname{Im}A)^\perp
	=
	\ker A^*.
	\]
	Следовательно,
	\[
	H=\overline{\operatorname{Im}A}\oplus\ker A^*.
	\]
\end{proof}

\begin{remark}
	Аналогично, применяя доказанное равенство к оператору \(A^*\), получаем
	\[
	H=\overline{\operatorname{Im}A^*}\oplus\ker A.
	\]
\end{remark}\newpage
	
	% XI. САМОСОПРЯЖЕННЫЕ ОПЕРАТОРЫ
	\section{Самосопряженные операторы}
	\hypertarget{bilet10}{}\setupnav{9}{11} \subsection{Свойства квадратичной формы \texorpdfstring{\((Ax,x)\)}{(Ax,x)} и собственных значений самосопряженного оператора \texorpdfstring{\(A\)}{A}}

\begin{definition}[Самосопряженный оператор]
	Пусть \(H\) --- гильбертово пространство над \(\mathbb{C}\), и
	\[
	A\in\mathcal{L}(H).
	\]
	Оператор \(A\) называется \textbf{самосопряженным}, если
	\[
	A^*=A.
	\]
	Эквивалентно,
	\[
	(Ax,y)=(x,Ay)
	\qquad
	\forall x,y\in H.
	\]
\end{definition}

\begin{definition}[Квадратичная форма самосопряженного оператора]
	Пусть \(A\in\mathcal{L}(H)\). С оператором \(A\) связывают квадратичную форму
	\[
	K:H\to\mathbb{C},
	\]
	заданную правилом
	\[
	K(x)=(Ax,x).
	\]
	Идея состоит в том, что вместо оператора \(A\) часто удобно изучать значения формы
	\[
	(Ax,x).
	\]
\end{definition}

\begin{theorem}[Свойства квадратичной формы и собственных значений]
	Пусть \(H\) --- гильбертово пространство над \(\mathbb{C}\), и
	\[
	A\in\mathcal{L}(H)
	\]
	--- самосопряженный оператор. Тогда:
	
	\begin{itemize}
		\item для любого \(x\in H\)
		\[
		(Ax,x)\in\mathbb{R};
		\]
		
		\item если \(\lambda\) --- собственное значение оператора \(A\), то
		\[
		\lambda\in\mathbb{R};
		\]
		
		\item если \(\lambda_1\neq\lambda_2\) --- собственные значения оператора \(A\), а
		\(e_1,e_2\) --- соответствующие им собственные векторы, то
		\[
		(e_1,e_2)=0.
		\]
	\end{itemize}
\end{theorem}

\begin{proof}
	Докажем три утверждения по отдельности.
	
	\textbf{1. Значения квадратичной формы вещественны.}
	
	Пусть
	\[
	x\in H.
	\]
	Так как \(A\) самосопряжен, имеем
	\[
	(Ax,x)=(x,Ax).
	\]
	По свойству скалярного произведения
	\[
	(x,Ax)=\overline{(Ax,x)}.
	\]
	Следовательно,
	\[
	(Ax,x)=\overline{(Ax,x)}.
	\]
	Число совпадает со своим комплексным сопряжением, значит
	\[
	(Ax,x)\in\mathbb{R}.
	\]
	
	\textbf{2. Собственные значения самосопряженного оператора вещественны.}
	
	Пусть \(\lambda\) --- собственное значение оператора \(A\). Тогда существует ненулевой вектор
	\[
	e\in H,
	\qquad
	e\neq0,
	\]
	такой что
	\[
	Ae=\lambda e.
	\]
	Рассмотрим значение квадратичной формы на \(e\):
	\[
	(Ae,e).
	\]
	С одной стороны, по первому пункту
	\[
	(Ae,e)\in\mathbb{R}.
	\]
	С другой стороны,
	\[
	(Ae,e)
	=
	(\lambda e,e)
	=
	\lambda(e,e)
	=
	\lambda\|e\|^2.
	\]
	Так как \(e\neq0\), то
	\[
	\|e\|^2>0.
	\]
	Значит,
	\[
	\lambda
	=
	\frac{(Ae,e)}{\|e\|^2}
	\in\mathbb{R}.
	\]
	
	\textbf{3. Собственные векторы разных собственных значений ортогональны.}
	
	Пусть
	\[
	Ae_1=\lambda_1 e_1,
	\qquad
	Ae_2=\lambda_2 e_2,
	\]
	где
	\[
	\lambda_1\neq\lambda_2.
	\]
	Тогда
	\[
	(Ae_1,e_2)
	=
	(\lambda_1 e_1,e_2)
	=
	\lambda_1(e_1,e_2).
	\]
	С другой стороны, так как \(A\) самосопряжен,
	\[
	(Ae_1,e_2)
	=
	(e_1,Ae_2)
	=
	(e_1,\lambda_2 e_2).
	\]
	По пункту 2 собственные значения самосопряженного оператора вещественны, поэтому
	\[
	(e_1,\lambda_2 e_2)
	=
	\lambda_2(e_1,e_2).
	\]
	Значит,
	\[
	\lambda_1(e_1,e_2)
	=
	\lambda_2(e_1,e_2).
	\]
	Отсюда
	\[
	(\lambda_1-\lambda_2)(e_1,e_2)=0.
	\]
	Так как
	\[
	\lambda_1\neq\lambda_2,
	\]
	получаем
	\[
	(e_1,e_2)=0.
	\]
\end{proof}

\begin{remark}[Обратное утверждение]
	Верно и обратное: если для любого \(x\in H\)
	\[
	(Ax,x)\in\mathbb{R},
	\]
	то оператор \(A\) самосопряжен.
	
	Это доказывается поляризацией: рассматривают значения формы на \(x+y\) и \(x+iy\), после чего
	из вещественности соответствующих выражений получают
	\[
	(Ax,y)=(x,Ay).
	\]
	Для основного билета это утверждение обычно можно оставить как дополнительное.
\end{remark}\newpage
	\hypertarget{bilet11}{}\setupnav{10}{12} \subsection{Разложение гильбертова пространства \texorpdfstring{\(H=\overline{\operatorname{Im}(A-\lambda I)}\oplus\ker(A-\lambda I)\)}{H = cl Im(A-lambda I) plus Ker(A-lambda I)}, где \texorpdfstring{\(A\)}{A} --- самосопряженный оператор}

\begin{definition}[Оператор \(A_\lambda\)]
	Пусть \(H\) --- гильбертово пространство над \(\mathbb{C}\),
	\[
	A\in\mathcal{L}(H),
	\qquad
	\lambda\in\mathbb{C}.
	\]
	Обозначим
	\[
	A_\lambda=A-\lambda I.
	\]
\end{definition}

\begin{theorem}
	Пусть \(H\) --- гильбертово пространство над \(\mathbb{C}\), и
	\[
	A\in\mathcal{L}(H)
	\]
	--- самосопряженный оператор. Тогда для любого \(\lambda\in\mathbb{C}\)
	\[
	H=\overline{\operatorname{Im}A_\lambda}\oplus_{\perp}\ker A_\lambda.
	\]
	То есть
	\[
	H
	=
	\overline{\operatorname{Im}(A-\lambda I)}
	\oplus_{\perp}
	\ker(A-\lambda I).
	\]
\end{theorem}

\begin{proof}
	Будем использовать разложение из билета 10.2:
	\[
	H=\overline{\operatorname{Im}B}\oplus_{\perp}\ker B^*
	\]
	для любого
	\[
	B\in\mathcal{L}(H).
	\]
	Применим его к оператору
	\[
	B=A_\lambda=A-\lambda I.
	\]
	Тогда
	\[
	H=\overline{\operatorname{Im}A_\lambda}\oplus_{\perp}\ker A_\lambda^*.
	\]
	Остаётся понять, почему в данном случае
	\[
	\ker A_\lambda^*=\ker A_\lambda.
	\]
	
	\textbf{1. Случай \(\lambda\in\mathbb{R}\).}
	
	Так как \(A\) самосопряжен, то
	\[
	A^*=A.
	\]
	Кроме того,
	\[
	(\lambda I)^*=\overline{\lambda}I.
	\]
	Поэтому
	\[
	A_\lambda^*
	=
	(A-\lambda I)^*
	=
	A^*-\overline{\lambda}I
	=
	A-\overline{\lambda}I.
	\]
	Если \(\lambda\in\mathbb{R}\), то
	\[
	\overline{\lambda}=\lambda.
	\]
	Значит,
	\[
	A_\lambda^*=A-\lambda I=A_\lambda.
	\]
	Следовательно,
	\[
	\ker A_\lambda^*=\ker A_\lambda.
	\]
	Подставляя это в разложение из 10.2, получаем
	\[
	H=\overline{\operatorname{Im}A_\lambda}\oplus_{\perp}\ker A_\lambda.
	\]
	
	\textbf{2. Случай \(\lambda\notin\mathbb{R}\).}
	
	По формуле выше
	\[
	A_\lambda^*=A-\overline{\lambda}I.
	\]
	Так как \(A\) самосопряжен, по билету 11.1 все собственные значения \(A\) вещественны.
	
	Но
	\[
	\lambda\notin\mathbb{R},
	\qquad
	\overline{\lambda}\notin\mathbb{R}.
	\]
	Значит, ни \(\lambda\), ни \(\overline{\lambda}\) не являются собственными значениями оператора \(A\).
	Следовательно,
	\[
	\ker(A-\lambda I)=\{0\},
	\]
	и
	\[
	\ker(A-\overline{\lambda}I)=\{0\}.
	\]
	То есть
	\[
	\ker A_\lambda=\{0\},
	\qquad
	\ker A_\lambda^*=\{0\}.
	\]
	
	Из разложения из 10.2 имеем
	\[
	H=\overline{\operatorname{Im}A_\lambda}\oplus_{\perp}\ker A_\lambda^*.
	\]
	Но \(\ker A_\lambda^*=\{0\}\), значит
	\[
	H=\overline{\operatorname{Im}A_\lambda}.
	\]
	Так как одновременно \(\ker A_\lambda=\{0\}\), можно записать
	\[
	H=\overline{\operatorname{Im}A_\lambda}\oplus_{\perp}\ker A_\lambda.
	\]
\end{proof}\newpage
	\hypertarget{bilet12}{}\setupnav{11}{13} \subsection{Критерий принадлежности числа спектру самосопряженного оператора. Вещественность спектра самосопряженного оператора}

\begin{definition}[Оператор \(A_\lambda\)]
	Пусть \(H\) --- гильбертово пространство над \(\mathbb{C}\),
	\[
	A\in\mathcal{L}(H),
	\qquad
	\lambda\in\mathbb{C}.
	\]
	Обозначим
	\[
	A_\lambda=A-\lambda I.
	\]
\end{definition}

\begin{remark}[Используемые факты]
	В доказательстве используем два уже доказанных факта.
	
	\begin{itemize}
		\item Если оператор ограничен снизу, то он инъективен и его образ замкнут.
		
		\item Для самосопряженного оператора \(A\) выполнено разложение
		\[
		H=\overline{\operatorname{Im}A_\lambda}\oplus_{\perp}\ker A_\lambda.
		\]
	\end{itemize}
\end{remark}

\begin{theorem}[Критерий принадлежности числа спектру]
	Пусть \(H\) --- гильбертово пространство над \(\mathbb{C}\), и
	\[
	A\in\mathcal{L}(H)
	\]
	--- самосопряженный оператор. Тогда:
	
	\begin{itemize}
		\item \[
		\lambda\in\rho(A)
		\Longleftrightarrow
		A_\lambda \text{ ограничен снизу},
		\]
		то есть
		\[
		\exists m>0
		\quad
		\forall x\in H
		\qquad
		\|A_\lambda x\|\geqslant m\|x\|.
		\]
		
		\item \[
		\lambda\in\sigma(A)
		\Longleftrightarrow
		\exists \{x_n\}\subset H:
		\|x_n\|=1,
		\quad
		A_\lambda x_n\to0.
		\]
	\end{itemize}
\end{theorem}

\begin{proof}
	Сначала докажем критерий через ограниченность снизу, а затем аккуратно перепишем его
	отрицание.
	
	\textbf{1. Если \(\lambda\in\rho(A)\), то \(A_\lambda\) ограничен снизу.}
	
	Пусть
	\[
	\lambda\in\rho(A).
	\]
	Тогда по определению резольвентного множества существует ограниченный обратный оператор
	\[
	A_\lambda^{-1}\in\mathcal{L}(H).
	\]
	По критерию из билета 8.1 отсюда следует, что \(A_\lambda\) ограничен снизу.
	
	То есть существует \(m>0\), такое что
	\[
	\|A_\lambda x\|\geqslant m\|x\|
	\qquad
	\forall x\in H.
	\]
	
	\textbf{2. Если \(A_\lambda\) ограничен снизу, то \(\lambda\in\rho(A)\).}
	
	Пусть существует \(m>0\), такое что
	\[
	\|A_\lambda x\|\geqslant m\|x\|
	\qquad
	\forall x\in H.
	\]
	Нужно доказать, что
	\[
	A_\lambda^{-1}\in\mathcal{L}(H).
	\]
	По билету 8.1 для этого достаточно показать, что
	\[
	\operatorname{Im}A_\lambda=H.
	\]
	
	\begin{itemize}
		\item \textbf{Ядро тривиально.}
		
		Если
		\[
		x\in\ker A_\lambda,
		\]
		то
		\[
		A_\lambda x=0.
		\]
		По ограниченности снизу
		\[
		0=\|A_\lambda x\|\geqslant m\|x\|.
		\]
		Так как \(m>0\), получаем
		\[
		x=0.
		\]
		Значит,
		\[
		\ker A_\lambda=\{0\}.
		\]
		
		\item \textbf{Образ плотен.}
		
		По разложению из билета 11.2
		\[
		H=\overline{\operatorname{Im}A_\lambda}\oplus_{\perp}\ker A_\lambda.
		\]
		Так как
		\[
		\ker A_\lambda=\{0\},
		\]
		получаем
		\[
		H=\overline{\operatorname{Im}A_\lambda}.
		\]
		
		\item \textbf{Образ замкнут.}
		
		Так как \(A_\lambda\) ограничен снизу, то по следствию из билета 8.1 его образ замкнут:
		\[
		\operatorname{Im}A_\lambda
		\text{ замкнут}.
		\]
		Но уже известно, что
		\[
		\overline{\operatorname{Im}A_\lambda}=H.
		\]
		Следовательно,
		\[
		\operatorname{Im}A_\lambda=H.
		\]
	\end{itemize}
	
	Итак, \(A_\lambda\) инъективен и сюръективен на \(H\), а также ограничен снизу. Поэтому по
	критерию из билета 8.1
	\[
	A_\lambda^{-1}\in\mathcal{L}(H).
	\]
	Значит,
	\[
	\lambda\in\rho(A).
	\]
	
	\textbf{3. Перепишем отрицание критерия.}
	
	Из первого пункта имеем
	\[
	\lambda\in\sigma(A)
	\Longleftrightarrow
	\lambda\notin\rho(A)
	\Longleftrightarrow
	A_\lambda \text{ не ограничен снизу}.
	\]
	Отрицание ограниченности снизу имеет вид:
	\[
	\forall m>0
	\quad
	\exists x\in H
	\qquad
	\|A_\lambda x\|<m\|x\|.
	\]
	Такой \(x\) автоматически ненулевой, иначе получилось бы \(0<0\).
	
	Выберем
	\[
	m=\frac{1}{n}.
	\]
	Тогда существует \(u_n\neq0\), такое что
	\[
	\|A_\lambda u_n\|<\frac{1}{n}\|u_n\|.
	\]
	Положим
	\[
	x_n=\frac{u_n}{\|u_n\|}.
	\]
	Тогда
	\[
	\|x_n\|=1
	\]
	и
	\[
	\|A_\lambda x_n\|
	=
	\frac{\|A_\lambda u_n\|}{\|u_n\|}
	<
	\frac{1}{n}.
	\]
	Следовательно,
	\[
	A_\lambda x_n\to0.
	\]
	
	Обратно, если существует последовательность \(\{x_n\}\), такая что
	\[
	\|x_n\|=1,
	\qquad
	A_\lambda x_n\to0,
	\]
	то \(A_\lambda\) не может быть ограничен снизу. Действительно, если бы существовало
	\(m>0\), такое что
	\[
	\|A_\lambda x\|\geqslant m\|x\|
	\qquad
	\forall x\in H,
	\]
	то для \(x=x_n\) получили бы
	\[
	\|A_\lambda x_n\|\geqslant m,
	\]
	что противоречит сходимости
	\[
	A_\lambda x_n\to0.
	\]
	Значит,
	\[
	\lambda\in\sigma(A).
	\]
\end{proof}

\begin{theorem}[Вещественность спектра самосопряженного оператора]
	Пусть \(H\) --- гильбертово пространство над \(\mathbb{C}\), и
	\[
	A\in\mathcal{L}(H)
	\]
	--- самосопряженный оператор. Тогда
	\[
	\sigma(A)\subset\mathbb{R}.
	\]
\end{theorem}

\begin{proof}
	Докажем равносильное утверждение:
	\[
	\lambda\notin\mathbb{R}
	\Longrightarrow
	\lambda\in\rho(A).
	\]
	
	\textbf{1. Представим \(\lambda\) и выделим самосопряженную часть.}
	
	Пусть
	\[
	\lambda=\mu+i\nu,
	\qquad
	\mu,\nu\in\mathbb{R},
	\qquad
	\nu\neq0.
	\]
	Тогда
	\[
	A_\lambda
	=
	A-\lambda I
	=
	A-\mu I-i\nu I.
	\]
	Обозначим
	\[
	A_\mu=A-\mu I.
	\]
	Так как \(A\) самосопряжен и \(\mu\in\mathbb{R}\), оператор \(A_\mu\) тоже самосопряжен.
	
	\textbf{2. Докажем оценку снизу для \(A_\lambda\).}
	
	Возьмём произвольный \(x\in H\). Тогда
	\[
	A_\lambda x
	=
	A_\mu x-i\nu x.
	\]
	Поэтому
	\[
	\|A_\lambda x\|^2
	=
	(A_\mu x-i\nu x,\ A_\mu x-i\nu x).
	\]
	Раскрывая скалярное произведение, получаем
	\[
	\|A_\lambda x\|^2
	=
	\|A_\mu x\|^2
	-i\nu(x,A_\mu x)
	+i\nu(A_\mu x,x)
	+|\nu|^2\|x\|^2.
	\]
	Так как \(A_\mu\) самосопряжен, по билету 11.1
	\[
	(A_\mu x,x)\in\mathbb{R}.
	\]
	Значит,
	\[
	(x,A_\mu x)=(A_\mu x,x),
	\]
	и смешанные слагаемые сокращаются:
	\[
	-i\nu(x,A_\mu x)
	+i\nu(A_\mu x,x)=0.
	\]
	Следовательно,
	\[
	\|A_\lambda x\|^2
	=
	\|A_\mu x\|^2
	+
	|\nu|^2\|x\|^2.
	\]
	Отсюда
	\[
	\|A_\lambda x\|^2
	\geqslant
	|\nu|^2\|x\|^2.
	\]
	Значит,
	\[
	\|A_\lambda x\|
	\geqslant
	|\nu|\|x\|.
	\]
	То есть \(A_\lambda\) ограничен снизу.
	
	\textbf{3. Применим критерий принадлежности спектру.}
	
	По доказанному критерию
	\[
	A_\lambda \text{ ограничен снизу}
	\Longrightarrow
	\lambda\in\rho(A).
	\]
	Следовательно,
	\[
	\lambda\notin\mathbb{R}
	\Longrightarrow
	\lambda\in\rho(A).
	\]
	То есть нерегулярных точек вне вещественной оси нет, поэтому
	\[
	\sigma(A)\subset\mathbb{R}.
	\]
\end{proof}
\newpage
	\hypertarget{bilet13}{}\setupnav{12}{14} \subsection{Теорема о спектре самосопряженного оператора: \texorpdfstring{\(\sigma(A)\subset[m_-,m_+]\), \(r(A)=\|A\|\)}{sigma(A) subset [m-,m+], r(A)=||A||}}

\begin{definition}[Числа \(m_-\) и \(m_+\)]
	Пусть \(H\) --- комплексное гильбертово пространство, и
	\[
	A\in\mathcal{L}(H)
	\]
	--- самосопряженный оператор.
	
	Положим
	\[
	m_-=\inf_{\|x\|=1}(Ax,x),
	\qquad
	m_+=\sup_{\|x\|=1}(Ax,x).
	\]
	Так как \(A\) самосопряжен, то
	\[
	(Ax,x)\in\mathbb{R}.
	\]
	Кроме того, при \(\|x\|=1\)
	\[
	|(Ax,x)|\leqslant \|A\|.
	\]
	Поэтому числа \(m_-\) и \(m_+\) корректно определены и конечны.
\end{definition}

\begin{remark}[Полускалярное произведение]
	Если \(B\in\mathcal{L}(H)\) самосопряжен и
	\[
	(Bx,x)\geqslant0
	\qquad
	\forall x\in H,
	\]
	то формула
	\[
	\langle x,y\rangle_B=(Bx,y)
	\]
	задаёт полускалярное произведение.
	
	Для него верно неравенство Коши--Буняковского:
	\[
	|\langle x,y\rangle_B|^2
	\leqslant
	\langle x,x\rangle_B\langle y,y\rangle_B.
	\]
\end{remark}

\begin{remark}[Используемые факты]
	В доказательстве используем уже доказанные факты:
	
	\begin{itemize}
		\item спектр самосопряженного оператора вещественен:
		\[
		\sigma(A)\subset\mathbb{R};
		\]
		
		\item критерий принадлежности спектру:
		\[
		\lambda\in\sigma(A)
		\Longleftrightarrow
		\exists x_n:\|x_n\|=1,\quad (A-\lambda I)x_n\to0;
		\]
		
		\item основная теорема о спектральном радиусе:
		\[
		r(A)=\lim_{n\to\infty}\sqrt[n]{\|A^n\|}.
		\]
	\end{itemize}
\end{remark}

\begin{theorem}[Спектр самосопряженного оператора]
	Пусть \(H\) --- комплексное гильбертово пространство, и
	\[
	A\in\mathcal{L}(H)
	\]
	--- самосопряженный оператор. Тогда
	\[
	\sigma(A)\subset[m_-,m_+],
	\]
	причём
	\[
	m_-,m_+\in\sigma(A).
	\]
	Кроме того,
	\[
	r(A)=\max\{|m_-|,|m_+|\}=\|A\|.
	\]
\end{theorem}

\begin{proof}
	Докажем утверждение по смысловым шагам.
	
	\textbf{1. Спектр лежит внутри отрезка \([m_-,m_+]\).}
	
	Так как \(A\) самосопряжен, то по предыдущему билету
	\[
	\sigma(A)\subset\mathbb{R}.
	\]
	Поэтому достаточно доказать, что все вещественные \(\lambda\), лежащие вне отрезка
	\([m_-,m_+]\), являются регулярными.
	
	\begin{itemize}
		\item \textbf{Случай \(\lambda>m_+\).}
		
		Пусть
		\[
		\lambda>m_+.
		\]
		Для любого \(x\in H\), \(x\neq0\), имеем
		\[
		\|A_\lambda x\|\cdot\|x\|
		\geqslant
		|(A_\lambda x,x)|.
		\]
		Но
		\[
		(A_\lambda x,x)
		=
		((A-\lambda I)x,x)
		=
		(Ax,x)-\lambda\|x\|^2.
		\]
		По определению \(m_+\)
		\[
		(Ax,x)\leqslant m_+\|x\|^2.
		\]
		Так как \(\lambda>m_+\), получаем
		\[
		|(Ax,x)-\lambda\|x\|^2|
		=
		\lambda\|x\|^2-(Ax,x)
		\geqslant
		(\lambda-m_+)\|x\|^2.
		\]
		Следовательно,
		\[
		\|A_\lambda x\|\cdot\|x\|
		\geqslant
		(\lambda-m_+)\|x\|^2.
		\]
		Сокращая на \(\|x\|\), получаем
		\[
		\|A_\lambda x\|
		\geqslant
		(\lambda-m_+)\|x\|.
		\]
		Значит, \(A_\lambda\) ограничен снизу, поэтому по критерию из билета 11.3
		\[
		\lambda\in\rho(A).
		\]
		
		\item \textbf{Случай \(\lambda<m_-\).}
		
		Пусть
		\[
		\lambda<m_-.
		\]
		Тогда по определению \(m_-\)
		\[
		(Ax,x)\geqslant m_-\|x\|^2.
		\]
		Следовательно,
		\[
		|(Ax,x)-\lambda\|x\|^2|
		=
		(Ax,x)-\lambda\|x\|^2
		\geqslant
		(m_- - \lambda)\|x\|^2.
		\]
		Повторяя оценку через Коши--Буняковского, получаем
		\[
		\|A_\lambda x\|
		\geqslant
		(m_- - \lambda)\|x\|.
		\]
		Значит, \(A_\lambda\) ограничен снизу, и снова
		\[
		\lambda\in\rho(A).
		\]
	\end{itemize}
	
	Итак, все точки вне отрезка \([m_-,m_+]\) лежат в резольвентном множестве. Значит,
	\[
	\sigma(A)\subset[m_-,m_+].
	\]
	
	\textbf{2. Докажем, что \(m_+\in\sigma(A)\).}
	
	По определению \(m_+\) как супремума нужно построить последовательность почти собственных
	векторов для \(m_+\), а потом применить критерий из билета 11.3.
	
	\begin{itemize}
		\item \textbf{Выберем максимизирующую последовательность.}
		
		По определению супремума существует последовательность \(\{x_n\}\subset H\), такая что
		\[
		\|x_n\|=1
		\]
		и
		\[
		(Ax_n,x_n)\to m_+.
		\]
		
		\item \textbf{Введём неотрицательный самосопряженный оператор.}
		
		Рассмотрим оператор
		\[
		B_+=m_+I-A.
		\]
		Он самосопряжен. Кроме того, для любого \(x\in H\)
		\[
		(B_+x,x)
		=
		m_+\|x\|^2-(Ax,x)
		\geqslant0.
		\]
		Значит, формула
		\[
		\langle u,v\rangle_{B_+}=(B_+u,v)
		\]
		задаёт полускалярное произведение.
		
		\item \textbf{Применим Коши--Буняковского для полускалярного произведения.}
		
		Для этого полускалярного произведения имеем
		\[
		|\langle u,v\rangle_{B_+}|^2
		\leqslant
		\langle u,u\rangle_{B_+}\langle v,v\rangle_{B_+}.
		\]
		Берём
		\[
		u=x_n,
		\qquad
		v=B_+x_n.
		\]
		Тогда
		\[
		\|B_+x_n\|^4
		=
		|(B_+x_n,B_+x_n)|^2
		=
		|\langle x_n,B_+x_n\rangle_{B_+}|^2.
		\]
		По неравенству Коши--Буняковского
		\[
		\|B_+x_n\|^4
		\leqslant
		\langle x_n,x_n\rangle_{B_+}
		\langle B_+x_n,B_+x_n\rangle_{B_+}.
		\]
		
		\item \textbf{Покажем, что правая часть стремится к нулю.}
		
		Первый множитель равен
		\[
		\langle x_n,x_n\rangle_{B_+}
		=
		(B_+x_n,x_n)
		=
		m_+-(Ax_n,x_n)
		\to0.
		\]
		Второй множитель ограничен:
		\[
		\langle B_+x_n,B_+x_n\rangle_{B_+}
		=
		(B_+^2x_n,B_+x_n)
		\leqslant
		\|B_+^2x_n\|\cdot\|B_+x_n\|
		\leqslant
		\|B_+\|^3.
		\]
		Следовательно,
		\[
		\|B_+x_n\|^4\to0,
		\]
		то есть
		\[
		B_+x_n\to0.
		\]
		
		\item \textbf{Применим критерий принадлежности спектру.}
		
		Так как
		\[
		B_+x_n=(m_+I-A)x_n=-(A-m_+I)x_n,
		\]
		получаем
		\[
		(A-m_+I)x_n\to0.
		\]
		При этом
		\[
		\|x_n\|=1.
		\]
		По критерию из билета 11.3
		\[
		m_+\in\sigma(A).
		\]
	\end{itemize}
	
	\textbf{3. Докажем, что \(m_-\in\sigma(A)\).}
	
	Доказательство полностью аналогично предыдущему пункту, только вместо супремума берём
	инфимум.
	
	\begin{itemize}
		\item \textbf{Выберем минимизирующую последовательность.}
		
		По определению \(m_-\) существует последовательность \(\{y_n\}\subset H\), такая что
		\[
		\|y_n\|=1
		\]
		и
		\[
		(Ay_n,y_n)\to m_-.
		\]
		
		\item \textbf{Введём неотрицательный самосопряженный оператор.}
		
		Рассмотрим
		\[
		B_-=A-m_-I.
		\]
		Тогда \(B_-\) самосопряжен и
		\[
		(B_-x,x)
		=
		(Ax,x)-m_-\|x\|^2
		\geqslant0.
		\]
		
		\item \textbf{Повторим оценку из предыдущего пункта.}
		
		Для полускалярного произведения
		\[
		\langle u,v\rangle_{B_-}=(B_-u,v)
		\]
		аналогично получаем
		\[
		\|B_-y_n\|^4
		\leqslant
		\langle y_n,y_n\rangle_{B_-}
		\langle B_-y_n,B_-y_n\rangle_{B_-}.
		\]
		Здесь
		\[
		\langle y_n,y_n\rangle_{B_-}
		=
		(B_-y_n,y_n)
		=
		(Ay_n,y_n)-m_-
		\to0,
		\]
		а второй множитель ограничен сверху константой, не зависящей от \(n\).
		
		Значит,
		\[
		B_-y_n\to0.
		\]
		
		\item \textbf{Применим критерий принадлежности спектру.}
		
		Так как
		\[
		B_-y_n=(A-m_-I)y_n,
		\]
		получаем
		\[
		(A-m_-I)y_n\to0.
		\]
		При этом
		\[
		\|y_n\|=1.
		\]
		По критерию из билета 11.3
		\[
		m_-\in\sigma(A).
		\]
	\end{itemize}
	
	\textbf{4. Найдём спектральный радиус через \(m_-\) и \(m_+\).}
	
	Мы уже доказали, что
	\[
	\sigma(A)\subset[m_-,m_+],
	\]
	и при этом
	\[
	m_-,m_+\in\sigma(A).
	\]
	Поэтому
	\[
	r(A)
	=
	\sup_{\lambda\in\sigma(A)}|\lambda|
	=
	\max\{|m_-|,|m_+|\}.
	\]
	
	\textbf{5. Докажем равенство \(r(A)=\|A\|\).}
	
	Используем общую формулу спектрального радиуса:
	\[
	r(A)=\lim_{n\to\infty}\sqrt[n]{\|A^n\|}.
	\]
	Нужно показать, что для самосопряженного оператора этот предел равен \(\|A\|\).
	
	\begin{itemize}
		\item \textbf{Докажем равенство \(\|A^2\|=\|A\|^2\).}
		
		Оценка
		\[
		\|A^2\|\leqslant\|A\|^2
		\]
		следует из мультипликативности нормы оператора.
		
		Докажем обратную оценку. Для любого \(x\in H\), \(\|x\|=1\), имеем
		\[
		\|Ax\|^2
		=
		(Ax,Ax)
		=
		(A^2x,x).
		\]
		По неравенству Коши--Буняковского
		\[
		(A^2x,x)
		\leqslant
		\|A^2x\|\cdot\|x\|
		\leqslant
		\|A^2\|.
		\]
		Значит,
		\[
		\|Ax\|^2\leqslant\|A^2\|.
		\]
		Переходя к супремуму по всем \(\|x\|=1\), получаем
		\[
		\|A\|^2\leqslant\|A^2\|.
		\]
		Следовательно,
		\[
		\|A^2\|=\|A\|^2.
		\]
		
		\item \textbf{Перейдём к степеням \(2^k\).}
		
		Так как \(A\) самосопряжен, то \(A^2\), \(A^4\), \(\ldots\), \(A^{2^k}\) тоже самосопряжены.
		Применяя доказанное равенство последовательно, получаем
		\[
		\|A^{2^k}\|=\|A\|^{2^k}.
		\]
		
		\item \textbf{Применим формулу спектрального радиуса.}
		
		По общей формуле
		\[
		r(A)=\lim_{n\to\infty}\sqrt[n]{\|A^n\|}.
		\]
		Тогда по подпоследовательности \(n=2^k\)
		\[
		r(A)
		=
		\lim_{k\to\infty}\sqrt[2^k]{\|A^{2^k}\|}
		=
		\lim_{k\to\infty}\sqrt[2^k]{\|A\|^{2^k}}
		=
		\|A\|.
		\]
	\end{itemize}
	
	Итак,
	\[
	r(A)=\max\{|m_-|,|m_+|\}=\|A\|.
	\]
\end{proof}\newpage
	
	% XII. КОМПАКТНЫЕ ОПЕРАТОРЫ
	\section{Компактные операторы}
	\hypertarget{bilet14}{}\setupnav{13}{15} \providecommand{\weakconv}{\rightharpoonup}

\subsection{Свойства компактных операторов}

\begin{definition}[Предкомпактное множество]
	Пусть \(E\) --- нормированное пространство. Множество
	\[
	M\subset E
	\]
	называется \textbf{предкомпактным}, если его замыкание
	\[
	\overline{M}
	\]
	компактно.
\end{definition}

\begin{definition}[Компактный оператор]
	Пусть \(E_1,E_2\) --- нормированные пространства и
	\[
	A\in\mathcal{L}(E_1,E_2).
	\]
	Оператор \(A\) называется \textbf{компактным}, если образ любого ограниченного множества
	в \(E_1\) является предкомпактным множеством в \(E_2\).
	
	Множество компактных операторов обозначают
	\[
	\mathcal{K}(E_1,E_2).
	\]
	Если \(E_1=E_2=E\), пишут
	\[
	\mathcal{K}(E).
	\]
\end{definition}

\begin{remark}[Предкомпактность и вполне ограниченность]
	В банаховом пространстве предкомпактность эквивалентна вполне ограниченности.
	
	То есть множество \(M\) предкомпактно тогда и только тогда, когда для любого
	\[
	\varepsilon>0
	\]
	существует конечная \(\varepsilon\)-сеть для \(M\).
\end{remark}

\begin{lemma}[Критерий компактности через единичный шар]
	Пусть \(E_1,E_2\) --- нормированные пространства и
	\[
	A\in\mathcal{L}(E_1,E_2).
	\]
	Тогда \(A\) компактен тогда и только тогда, когда множество
	\[
	A\bigl(\overline{B}_1(0)\bigr)
	\]
	предкомпактно в \(E_2\).
\end{lemma}

\begin{proof}
	Докажем обе стороны.
	
	\textbf{1. Если \(A\) компактен, то \(A(\overline{B}_1(0))\) предкомпактно.}
	
	Единичный шар
	\[
	\overline{B}_1(0)
	\]
	ограничен. Поэтому по определению компактного оператора его образ
	\[
	A\bigl(\overline{B}_1(0)\bigr)
	\]
	предкомпактен.
	
	\textbf{2. Если \(A(\overline{B}_1(0))\) предкомпактно, то \(A\) компактен.}
	
	Пусть
	\[
	S\subset E_1
	\]
	--- произвольное ограниченное множество. Тогда существует \(R>0\), такое что
	\[
	S\subset R\overline{B}_1(0).
	\]
	Так как \(A\) линеен,
	\[
	A(S)\subset A\bigl(R\overline{B}_1(0)\bigr)
	=
	R A\bigl(\overline{B}_1(0)\bigr).
	\]
	Если \(A(\overline{B}_1(0))\) предкомпактно, то и его растяжение
	\[
	R A\bigl(\overline{B}_1(0)\bigr)
	\]
	предкомпактно. Следовательно, \(A(S)\) как подмножество предкомпактного множества тоже
	предкомпактно.
	
	Значит, \(A\) компактен.
\end{proof}

\begin{remark}[Последовательностный критерий компактности]
	Оператор
	\[
	A\in\mathcal{L}(E_1,E_2)
	\]
	компактен тогда и только тогда, когда для любой ограниченной последовательности
	\[
	\{x_n\}\subset E_1
	\]
	последовательность
	\[
	\{Ax_n\}\subset E_2
	\]
	имеет сходящуюся подпоследовательность.
\end{remark}

\begin{theorem}[Конечномерный случай]
	Пусть \(E_1,E_2\) --- нормированные пространства и
	\[
	A\in\mathcal{L}(E_1,E_2).
	\]
	Если
	\[
	\dim E_1<\infty
	\]
	или
	\[
	\dim E_2<\infty,
	\]
	то оператор \(A\) компактен.
\end{theorem}

\begin{proof}
	\textbf{1. Если \(\dim E_1<\infty\).}
	
	Тогда образ оператора
	\[
	\operatorname{Im}A
	\]
	тоже конечномерен, потому что
	\[
	\dim\operatorname{Im}A\leqslant \dim E_1.
	\]
	Если \(S\subset E_1\) ограничено, то \(A(S)\) ограничено в конечномерном пространстве
	\(\operatorname{Im}A\). В конечномерном пространстве всякое ограниченное множество
	предкомпактно. Следовательно, \(A(S)\) предкомпактно.
	
	Значит, \(A\) компактен.
	
	\textbf{2. Если \(\dim E_2<\infty\).}
	
	Пусть \(S\subset E_1\) ограничено. Так как \(A\) ограничен,
	\[
	A(S)
	\]
	тоже ограничено. Но \(A(S)\subset E_2\), а \(E_2\) конечномерно. Поэтому \(A(S)\)
	предкомпактно.
	
	Значит, \(A\) компактен.
\end{proof}

\begin{theorem}[Тождественный оператор не компактен в бесконечномерном пространстве]
	Пусть \(E\) --- бесконечномерное нормированное пространство. Тогда
	\[
	I\notin\mathcal{K}(E).
	\]
\end{theorem}

\begin{proof}
	Предположим противное:
	\[
	I\in\mathcal{K}(E).
	\]
	Тогда образ единичного шара под действием \(I\) предкомпактен:
	\[
	I\bigl(\overline{B}_1(0)\bigr)=\overline{B}_1(0).
	\]
	Значит, единичный шар был бы предкомпактен.
	
	Но по теореме Рисса в бесконечномерном нормированном пространстве единичный шар не
	является компактным и даже не является вполне ограниченным.
	
	Получили противоречие. Следовательно,
	\[
	I\notin\mathcal{K}(E).
	\]
\end{proof}

\begin{theorem}[Компактные операторы образуют двусторонний идеал]
	Пусть \(E\) --- нормированное пространство. Тогда
	\[
	\mathcal{K}(E)
	\]
	является двусторонним идеалом в \(\mathcal{L}(E)\).
	
	То есть если
	\[
	A\in\mathcal{K}(E),
	\qquad
	B\in\mathcal{L}(E),
	\]
	то
	\[
	AB\in\mathcal{K}(E),
	\qquad
	BA\in\mathcal{K}(E).
	\]
\end{theorem}

\begin{proof}
	Докажем компактность обоих произведений.
	
	\textbf{1. Компактность \(AB\).}
	
	Пусть
	\[
	S\subset E
	\]
	--- ограниченное множество. Так как \(B\in\mathcal{L}(E)\), оператор \(B\) ограничен. Поэтому
	\[
	B(S)
	\]
	тоже ограничено.
	
	Так как \(A\) компактен, образ ограниченного множества \(B(S)\) под действием \(A\)
	предкомпактен:
	\[
	A(B(S))
	\text{ предкомпактно}.
	\]
	Но
	\[
	A(B(S))=(AB)(S).
	\]
	Значит, \(AB\) компактен.
	
	\textbf{2. Компактность \(BA\).}
	
	Пусть
	\[
	S\subset E
	\]
	--- ограниченное множество. Так как \(A\) компактен,
	\[
	A(S)
	\]
	предкомпактно.
	
	Нужно доказать, что
	\[
	B(A(S))
	\]
	предкомпактно.
	
	\begin{itemize}
		\item \textbf{Через компактное замыкание.}
		
		Так как \(A(S)\) предкомпактно, множество
		\[
		\overline{A(S)}
		\]
		компактно.
		
		Оператор \(B\) ограничен, значит непрерывен. Поэтому образ компакта при \(B\)
		компактен:
		\[
		B\bigl(\overline{A(S)}\bigr)
		\text{ компактен}.
		\]
		
		\item \textbf{Переход к нужному множеству.}
		
		Имеем
		\[
		B(A(S))\subset B\bigl(\overline{A(S)}\bigr).
		\]
		Следовательно, \(B(A(S))\) является подмножеством компактного множества, значит
		предкомпактно.
	\end{itemize}
	
	То есть
	\[
	(BA)(S)
	\]
	предкомпактно для любого ограниченного \(S\). Следовательно,
	\[
	BA\in\mathcal{K}(E).
	\]
\end{proof}

\begin{remark}[Линейность класса компактных операторов]
	Если
	\[
	A_1,A_2\in\mathcal{K}(E_1,E_2),
	\qquad
	\alpha,\beta\in\mathbb{K},
	\]
	то
	\[
	\alpha A_1+\beta A_2\in\mathcal{K}(E_1,E_2).
	\]
	Это следует из того, что сумма двух предкомпактных множеств предкомпактна, а умножение
	предкомпактного множества на число сохраняет предкомпактность.
\end{remark}

\begin{theorem}[Замкнутость класса компактных операторов по операторной норме]
	Пусть \(E_1\) --- нормированное пространство, \(E_2\) --- банахово пространство,
	\[
	A_n\in\mathcal{K}(E_1,E_2),
	\qquad
	A_n\to A
	\]
	по операторной норме:
	\[
	\|A_n-A\|\to0.
	\]
	Тогда
	\[
	A\in\mathcal{K}(E_1,E_2).
	\]
\end{theorem}

\begin{proof}
	По критерию через единичный шар достаточно доказать, что
	\[
	A\bigl(\overline{B}_1(0)\bigr)
	\]
	предкомпактно.
	
	Так как \(E_2\) банахово, достаточно доказать, что множество
	\[
	A\bigl(\overline{B}_1(0)\bigr)
	\]
	вполне ограничено.
	
	\textbf{1. Приблизим \(A\) компактным оператором.}
	
	Пусть
	\[
	\varepsilon>0.
	\]
	Так как
	\[
	\|A_n-A\|\to0,
	\]
	существует номер \(N\), такой что
	\[
	\|A_N-A\|<\frac{\varepsilon}{2}.
	\]
	Тогда для любого \(x\in\overline{B}_1(0)\)
	\[
	\|A_Nx-Ax\|
	\leqslant
	\|A_N-A\|\cdot\|x\|
	<
	\frac{\varepsilon}{2}.
	\]
	
	\textbf{2. Возьмём конечную сеть для \(A_N(\overline{B}_1(0))\).}
	
	Оператор \(A_N\) компактен, значит множество
	\[
	A_N\bigl(\overline{B}_1(0)\bigr)
	\]
	предкомпактно. Так как \(E_2\) банахово, оно вполне ограничено.
	
	Следовательно, существует конечная \(\varepsilon/2\)-сеть
	\[
	y_1,\ldots,y_m
	\]
	для множества
	\[
	A_N\bigl(\overline{B}_1(0)\bigr).
	\]
	То есть для любого \(x\in\overline{B}_1(0)\) найдётся номер \(j\in\{1,\ldots,m\}\), такой что
	\[
	\|A_Nx-y_j\|<\frac{\varepsilon}{2}.
	\]
	
	\textbf{3. Получим конечную сеть для \(A(\overline{B}_1(0))\).}
	
	Для такого \(x\) и соответствующего \(j\) имеем
	\[
	\|Ax-y_j\|
	\leqslant
	\|Ax-A_Nx\|+\|A_Nx-y_j\|
	<
	\varepsilon.
	\]
	Значит,
	\[
	y_1,\ldots,y_m
	\]
	является конечной \(\varepsilon\)-сетью для множества
	\[
	A\bigl(\overline{B}_1(0)\bigr).
	\]
	
	Так как \(\varepsilon>0\) произвольно, множество \(A(\overline{B}_1(0))\) вполне ограничено.
	Следовательно, оно предкомпактно.
	
	Значит,
	\[
	A\in\mathcal{K}(E_1,E_2).
	\]
\end{proof}

\begin{definition}[Вполне непрерывный оператор]
	Оператор \(A\in\mathcal{L}(E_1,E_2)\) называется \textbf{вполне непрерывным}, если из слабой
	сходимости
	\[
	x_n\weakconv x
	\]
	следует сильная сходимость образов:
	\[
	\|Ax_n-Ax\|\to0.
	\]
\end{definition}

\begin{theorem}[Компактный оператор улучшает слабую сходимость]
	Пусть \(E_1,E_2\) --- нормированные пространства и
	\[
	A\in\mathcal{K}(E_1,E_2).
	\]
	Если
	\[
	x_n\weakconv x
	\]
	в \(E_1\), то
	\[
	\|Ax_n-Ax\|\to0.
	\]
\end{theorem}

\begin{proof}
	Обозначим
	\[
	y_n=Ax_n.
	\]
	Нужно доказать, что
	\[
	y_n\to Ax
	\]
	по норме.
	
	\textbf{1. Последовательность образов предкомпактна.}
	
	Слабо сходящаяся последовательность ограничена. Поэтому множество
	\[
	\{x_n:n\in\mathbb{N}\}\cup\{x\}
	\]
	ограничено в \(E_1\).
	
	Так как \(A\) компактен, множество его образов
	\[
	\{Ax_n:n\in\mathbb{N}\}\cup\{Ax\}
	\]
	предкомпактно в \(E_2\).
	
	\textbf{2. Образы слабо сходятся к \(Ax\).}
	
	Так как \(A\) --- линейный ограниченный оператор, он сохраняет слабую сходимость:
	\[
	x_n\weakconv x
	\quad\Longrightarrow\quad
	Ax_n\weakconv Ax.
	\]
	То есть
	\[
	y_n\weakconv Ax.
	\]
	
	\textbf{3. Предкомпактность плюс слабая сходимость дают сильную сходимость.}
	
	Докажем от противного. Пусть
	\[
	y_n\not\to Ax
	\]
	по норме. Тогда существует \(\varepsilon_0>0\) и подпоследовательность \(\{y_{n_k}\}\), такая что
	\[
	\|y_{n_k}-Ax\|\geqslant\varepsilon_0
	\qquad
	\forall k.
	\]
	Но множество \(\{y_n\}\) предкомпактно, поэтому из \(\{y_{n_k}\}\) можно выделить
	подпоследовательность \(\{y_{n_{k_j}}\}\), сходящуюся по норме к некоторому \(y\in E_2\):
	\[
	y_{n_{k_j}}\to y.
	\]
	Сходимость по норме влечёт слабую сходимость:
	\[
	y_{n_{k_j}}\weakconv y.
	\]
	С другой стороны, так как вся последовательность \(\{y_n\}\) слабо сходится к \(Ax\), то
	\[
	y_{n_{k_j}}\weakconv Ax.
	\]
	Слабый предел единственен. Значит,
	\[
	y=Ax.
	\]
	То есть
	\[
	y_{n_{k_j}}\to Ax
	\]
	по норме, что противоречит условию
	\[
	\|y_{n_k}-Ax\|\geqslant\varepsilon_0.
	\]
	
	Следовательно,
	\[
	Ax_n\to Ax
	\]
	по норме.
\end{proof}

\begin{remark}[Обратное утверждение]
	Если пространство \(E_1\) рефлексивно, то компактность оператора эквивалентна полной
	непрерывности:
	\[
	A\in\mathcal{K}(E_1,E_2)
	\quad\Longleftrightarrow\quad
	x_n\weakconv x \Rightarrow \|Ax_n-Ax\|\to0.
	\]
	В общем случае обратное утверждение требует дополнительных условий.
\end{remark}
\newpage
	\hypertarget{bilet15}{}\setupnav{14}{16} \subsection{Свойства собственных значений компактного оператора}

\begin{definition}[Собственное значение]
	Пусть \(E\) --- комплексное банахово пространство и
	\[
	A\in\mathcal{L}(E).
	\]
	Число \(\lambda\in\mathbb{C}\) называется \textbf{собственным значением} оператора \(A\), если
	существует ненулевой вектор \(x\in E\), такой что
	\[
	Ax=\lambda x.
	\]
	Вектор \(x\) называется \textbf{собственным вектором}.
\end{definition}

\begin{remark}[Почему ноль особый]
	Для компактного оператора ненулевые собственные значения ведут себя почти как в конечномерной
	линейной алгебре: вне любой окрестности нуля их может быть только конечное число.
	
	Но \(0\) может быть собственным значением бесконечной кратности.
\end{remark}

\begin{example}
	Пусть оператор \(A:l_2\to l_2\) задаётся диагональной матрицей
	\[
	A=
	\begin{pmatrix}
		1 & 0 & 0 & \cdots\\
		0 & \frac12 & 0 & \cdots\\
		0 & 0 & \frac13 & \cdots\\
		\vdots & \vdots & \vdots & \ddots
	\end{pmatrix}.
	\]
	То есть
	\[
	A(x_1,x_2,x_3,\ldots)
	=
	\left(x_1,\frac{x_2}{2},\frac{x_3}{3},\ldots\right).
	\]
	Этот оператор компактен.
	
	Его ненулевые собственные значения:
	\[
	1,\frac12,\frac13,\ldots
	\]
	и они стремятся к нулю. При этом \(0\) является собственным значением бесконечной кратности
	для подходящих компактных операторов, поэтому \(0\) нельзя рассматривать так же, как
	ненулевые собственные значения.
\end{example}

\begin{theorem}[Собственные значения компактного самосопряженного оператора вне нуля]
	Пусть \(H\) --- гильбертово пространство и
	\[
	A\in\mathcal K(H)
	\]
	--- самосопряженный оператор. Тогда для любого
	\[
	\delta>0
	\]
	существует лишь конечное число собственных значений \(\lambda\), для которых
	\[
	|\lambda|\geq\delta.
	\]
\end{theorem}

\begin{proof}
	\textbf{Хотим получить противоречие с компактностью \(A\).}
	
	Предположим, что существуют различные собственные значения
	\[
	\lambda_1,\lambda_2,\ldots,
	\qquad
	|\lambda_n|\geq\delta.
	\]
	Выберем соответствующие нормированные собственные векторы
	\[
	Ae_n=\lambda_n e_n,
	\qquad
	\|e_n\|=1.
	\]
	
	\textbf{Хотим показать, что образы \(Ae_n\) не могут иметь сходящейся подпоследовательности.}
	
	Так как \(A\) самосопряжен, собственные векторы различных собственных значений
	ортогональны:
	\[
	e_n\perp e_m,
	\qquad n\neq m.
	\]
	
	Поэтому
	\[
	\begin{aligned}
		\|Ae_n-Ae_m\|^2
		&=
		\|\lambda_n e_n-\lambda_m e_m\|^2
		\\
		&=
		|\lambda_n|^2+|\lambda_m|^2
		\\
		&\geq
		2\delta^2.
	\end{aligned}
	\]
	Следовательно,
	\[
	\|Ae_n-Ae_m\|
	\geq
	\sqrt2\,\delta,
	\qquad n\neq m.
	\]
	
	Значит, последовательность
	\[
	\{Ae_n\}
	\]
	не содержит фундаментальной, а значит и сходящейся подпоследовательности.
	
	Но
	\[
	\|e_n\|=1,
	\]
	то есть \(\{e_n\}\) ограничена. Так как \(A\) компактен,
	последовательность \(\{Ae_n\}\) должна иметь сходящуюся подпоследовательность.
	
	Получили противоречие.
	
	Следовательно, для любого \(\delta>0\) существует лишь конечное число
	собственных значений, удовлетворяющих
	\[
	|\lambda|\geq\delta.
	\]
\end{proof}

\begin{remark}[Общий случай]
	В общем банаховом случае это утверждение доказывается похожей идеей, но вместо
	ортогональности используют теорему Рисса о почти перпендикуляре.
	
	На экзамене обычно достаточно понимать смысл: если бы ненулевых собственных значений
	вдали от нуля было бесконечно много, компактность дала бы сходящуюся подпоследовательность
	образов, но собственные векторы можно выбрать так, что образы остаются отделёнными друг от друга.
\end{remark}
\newpage
	\hypertarget{bilet16}{}\setupnav{15}{17} \subsection{Теорема Фредгольма для компактных самосопряженных операторов}

\begin{definition}[Оператор \(A_\lambda\)]
	Пусть
	\[
	A\in\mathcal L(H),
	\qquad
	\lambda\in\mathbb C.
	\]
	Обозначим
	\[
	A_\lambda=A-\lambda I.
	\]
\end{definition}

\begin{lemma}
	Пусть
	\[
	A\in\mathcal K(H)
	\]
	--- самосопряженный оператор и
	\[
	\lambda\neq0.
	\]
	Тогда на подпространстве
	\[
	M=(\ker A_\lambda)^\perp
	\]
	оператор \(A_\lambda\) ограничен снизу:
	\[
	\exists c>0
	\qquad
	\|A_\lambda x\|\geq c\|x\|,
	\qquad x\in M.
	\]
\end{lemma}

\begin{proof}
	\textbf{Хотим получить противоречие с компактностью \(A\).}
	
	Предположим, что \(A_\lambda\) не ограничен снизу на \(M\).
	Тогда существуют
	\[
	x_n\in M,
	\qquad
	\|x_n\|=1,
	\]
	такие что
	\[
	A_\lambda x_n\to0.
	\]
	
	Так как \(A\) компактен, из ограниченной последовательности \(\{x_n\}\)
	можно выбрать подпоследовательность, для которой
	\[
	Ax_{n_k}\to y.
	\]
	
	Но
	\[
	\lambda x_{n_k}
	=
	Ax_{n_k}-A_\lambda x_{n_k}
	\to y.
	\]
	Так как \(\lambda\neq0\),
	\[
	x_{n_k}\to x:=\frac{y}{\lambda}.
	\]
	
	Подпространство \(M\) замкнуто, поэтому
	\[
	x\in M,
	\qquad
	\|x\|=1.
	\]
	
	Переходя к пределу в
	\[
	A_\lambda x_{n_k}\to0,
	\]
	получаем
	\[
	A_\lambda x=0.
	\]
	Следовательно,
	\[
	x\in\ker A_\lambda.
	\]
	
	Таким образом,
	\[
	x\in M\cap\ker A_\lambda=\{0\},
	\]
	что противоречит
	\[
	\|x\|=1.
	\]
\end{proof}

\begin{lemma}
	При тех же условиях
	\[
	\operatorname{Im}A_\lambda
	\]
	замкнут.
\end{lemma}

\begin{proof}
	По предыдущей лемме оператор
	\[
	A_\lambda|_{(\ker A_\lambda)^\perp}
	\]
	ограничен снизу. Так как \((\ker A_\lambda)^\perp\) --- банахово пространство,
	его образ замкнут. Следовательно,
	\[
	\operatorname{Im}A_\lambda
	\]
	замкнут.
\end{proof}

\begin{theorem}[Фредгольма]
	Пусть
	\[
	A\in\mathcal K(H)
	\]
	--- компактный самосопряженный оператор и
	\[
	\lambda\neq0.
	\]
	Тогда
	\[
	H
	=
	\operatorname{Im}A_\lambda
	\oplus_{\perp}
	\ker A_\lambda.
	\]
\end{theorem}

\begin{proof}
	Так как \(A_\lambda\) самосопряжен, имеем стандартное разложение
	\[
	H
	=
	\overline{\operatorname{Im}A_\lambda}
	\oplus_{\perp}
	\ker A_\lambda.
	\]
	
	По предыдущей лемме
	\[
	\operatorname{Im}A_\lambda
	\]
	замкнут, поэтому
	\[
	\overline{\operatorname{Im}A_\lambda}
	=
	\operatorname{Im}A_\lambda.
	\]
	
	Следовательно,
	\[
	\boxed{
		H
		=
		\operatorname{Im}A_\lambda
		\oplus_{\perp}
		\ker A_\lambda.
	}
	\]
\end{proof}\newpage
	\hypertarget{bilet17}{}\setupnav{16}{18} \subsection{Теорема Гильберта--Шмидта}

\begin{definition}[Собственное подпространство]
	Пусть \(H\) --- комплексное гильбертово пространство,
	\[
	A\in\mathcal{L}(H),
	\qquad
	\lambda\in\mathbb{C}.
	\]
	Собственным подпространством, соответствующим \(\lambda\), называется
	\[
	H_\lambda=\ker(A-\lambda I).
	\]
	Если \(x\in H_\lambda\), \(x\neq0\), то
	\[
	Ax=\lambda x.
	\]
\end{definition}

\begin{lemma}[Наличие ненулевого собственного значения]
	Пусть \(H\) --- комплексное гильбертово пространство,
	\[
	A\in\mathcal{K}(H)
	\]
	--- компактный самосопряженный оператор. Если
	\[
	A\neq0,
	\]
	то у \(A\) существует ненулевое собственное значение.
\end{lemma}

\begin{proof}
	Так как \(A\) самосопряжен, по теореме о спектре самосопряженного оператора
	\[
	r(A)=\|A\|.
	\]
	Поскольку
	\[
	A\neq0,
	\]
	имеем
	\[
	\|A\|\neq0.
	\]
	Значит,
	\[
	r(A)>0.
	\]
	
	По той же теореме
	\[
	r(A)=\max\{|m_-|,|m_+|\},
	\]
	причём
	\[
	m_-,m_+\in\sigma(A).
	\]
	Следовательно, хотя бы одно из чисел \(m_-\), \(m_+\) отлично от нуля. Обозначим его через
	\[
	\lambda.
	\]
	Тогда
	\[
	\lambda\in\sigma(A),
	\qquad
	\lambda\neq0.
	\]
	
	По теореме о компактных самосопряженных операторах всякая ненулевая точка спектра является
	собственным значением. Поэтому \(\lambda\) --- ненулевое собственное значение оператора \(A\).
\end{proof}

\begin{theorem}[Гильберта--Шмидта]
	Пусть \(H\) --- сепарабельное комплексное гильбертово пространство, и
	\[
	A\in\mathcal{K}(H)
	\]
	--- компактный самосопряженный оператор. Тогда в \(H\) существует ортонормированный базис,
	состоящий из собственных векторов оператора \(A\).
\end{theorem}

\begin{proof}
	Докажем, что собственных векторов оператора \(A\) достаточно, чтобы породить всё пространство.
	
	\textbf{1. Отдельно рассмотрим тривиальный случай.}
	
	Если
	\[
	A=0,
	\]
	то любой ненулевой вектор является собственным вектором с собственным значением \(0\).
	Так как \(H\) сепарабельно, в \(H\) существует ортонормированный базис. Он автоматически
	состоит из собственных векторов оператора \(A\).
	
	Далее считаем, что
	\[
	A\neq0.
	\]
	
	\textbf{2. Соберём собственные векторы, отвечающие ненулевым собственным значениям.}
	
	По свойствам собственных значений компактного оператора ненулевых собственных значений
	не более чем счётно. Обозначим их через
	\[
	\lambda_1,\lambda_2,\ldots
	\]
	с учётом возможной конечности этого набора.
	
	Для каждого \(\lambda_j\neq0\) собственное подпространство
	\[
	H_{\lambda_j}=\ker(A-\lambda_j I)
	\]
	конечномерно.
	
	\begin{itemize}
		\item \textbf{Выберем базисы в собственных подпространствах.}
		
		В каждом \(H_{\lambda_j}\) выберем ортонормированный базис:
		\[
		e^{(j)}_1,\ldots,e^{(j)}_{m_j}.
		\]
		
		\item \textbf{Объединим эти базисы.}
		
		Все выбранные векторы являются собственными векторами оператора \(A\). Кроме того,
		собственные векторы, отвечающие различным собственным значениям самосопряженного
		оператора, ортогональны.
		
		Значит, система всех выбранных векторов ортонормирована.
	\end{itemize}
	
	Обозначим через \(M\) замкнутое подпространство, натянутое на все эти векторы:
	\[
	M=
	\overline{\operatorname{span}}
	\{e^{(j)}_k:\lambda_j\neq0,\ 1\leqslant k\leqslant m_j\}.
	\]
	
	\textbf{3. Покажем, что на \(M^\perp\) оператор \(A\) равен нулю.}
	
	Сначала заметим, что \(M\) инвариантно относительно \(A\). Действительно, на каждом собственном
	векторе
	\[
	Ae^{(j)}_k=\lambda_j e^{(j)}_k.
	\]
	Значит,
	\[
	A(M)\subset M.
	\]
	
	Так как \(A\) самосопряжен, из инвариантности \(M\) следует инвариантность \(M^\perp\):
	\[
	A(M^\perp)\subset M^\perp.
	\]
	
	Рассмотрим сужение
	\[
	\widetilde A=A|_{M^\perp}:M^\perp\to M^\perp.
	\]
	Оператор \(\widetilde A\) компактен и самосопряжен.
	
	\begin{itemize}
		\item \textbf{Предположим, что \(\widetilde A\neq0\).}
		
		Тогда по лемме у \(\widetilde A\) существует ненулевое собственное значение \(\mu\neq0\).
		Значит, существует ненулевой вектор
		\[
		u\in M^\perp,
		\]
		такой что
		\[
		\widetilde A u=\mu u.
		\]
		
		\item \textbf{Получим противоречие.}
		
		Так как \(\widetilde A\) --- это сужение \(A\), имеем
		\[
		Au=\mu u.
		\]
		Значит, \(u\) --- собственный вектор оператора \(A\), отвечающий ненулевому собственному
		значению \(\mu\).
		
		По построению все собственные векторы, отвечающие ненулевым собственным значениям,
		лежат в \(M\). Следовательно,
		\[
		u\in M.
		\]
		Но одновременно
		\[
		u\in M^\perp.
		\]
		Значит,
		\[
		u\in M\cap M^\perp=\{0\},
		\]
		что противоречит выбору \(u\neq0\).
	\end{itemize}
	
	Следовательно,
	\[
	\widetilde A=0.
	\]
	То есть
	\[
	Ax=0
	\qquad
	\forall x\in M^\perp.
	\]
	Значит,
	\[
	M^\perp\subset\ker A.
	\]
	
	\textbf{4. Добавим собственные векторы для собственного значения \(0\).}
	
	С другой стороны, если
	\[
	x\in\ker A,
	\]
	то
	\[
	Ax=0.
	\]
	Для любого собственного вектора \(e\), отвечающего ненулевому собственному значению \(\lambda\),
	по ортогональности собственных векторов самосопряженного оператора получаем
	\[
	(x,e)=0.
	\]
	Значит,
	\[
	\ker A\subset M^\perp.
	\]
	Итак,
	\[
	M^\perp=\ker A.
	\]
	
	Так как \(\ker A\) --- замкнутое подпространство сепарабельного гильбертова пространства,
	в нём существует ортонормированный базис. Обозначим его через
	\[
	f_1,f_2,\ldots
	\]
	с учётом возможной конечности или пустоты этого набора.
	
	Все эти векторы являются собственными векторами оператора \(A\), соответствующими собственному
	значению \(0\).
	
	\textbf{5. Получим ортонормированный базис всего пространства.}
	
	Мы получили ортогональное разложение
	\[
	H=M\oplus M^\perp
	=
	M\oplus\ker A.
	\]
	В \(M\) уже выбран ортонормированный базис из собственных векторов, отвечающих ненулевым
	собственным значениям. В \(\ker A\) выбран ортонормированный базис из собственных векторов,
	отвечающих собственному значению \(0\).
	
	Объединяя эти две системы, получаем ортонормированную систему
	\[
	\{e^{(j)}_k\}\cup\{f_s\}.
	\]
	Она полна в \(H\), потому что порождает
	\[
	M\oplus\ker A=H.
	\]
	Следовательно, это ортонормированный базис \(H\), состоящий из собственных векторов оператора
	\(A\).
\end{proof}\newpage
	
	% XIII. ЭЛЕМЕНТЫ НЕЛИНЕЙНОГО АНАЛИЗА
	\section{Элементы нелинейного анализа}
	\hypertarget{bilet18}{}\setupnav{17}{19} % --- Содержимое файла materials/13_1.tex ---

\subsection{Производная Фреше, производная Гато. Формула конечных приращений.}

\begin{remark}
	TODO: сюда дописывается содержимое билета в том же стиле, что и в конспекте по функциональному анализу--1.
	Сначала фиксируются нужные определения и формулировки теорем, затем подробно расписываются доказательства без пропуска логических переходов.
\end{remark}
\newpage
	
	% XIV. ПРЕОБРАЗОВАНИЕ ФУРЬЕ И СВЕРТКА В ПРОСТРАНСТВАХ $L_1(\MATHBB{R})$ И $L_2(\MATHBB{R})$
	\section{Преобразование Фурье и свертка в пространствах $L_1(\mathbb{R})$ и $L_2(\mathbb{R})$}
	\hypertarget{bilet19}{}\setupnav{18}{20} \subsection{Определения и основные свойства. Формула умножения. Преобразование Фурье свертки.}

\begin{definition}[Преобразование Фурье в \(L_1(\mathbb R)\)]
	Пусть
	\[
	f\in L_1(\mathbb R).
	\]
	Преобразованием Фурье функции \(f\) называется функция
	\[
	\widehat f(y)
	=
	\int_{\mathbb R} f(x)e^{-ixy}\,dx.
	\]
\end{definition}

\begin{theorem}[Основные свойства преобразования Фурье]
	Для любого \(f\in L_1(\mathbb R)\):
	\begin{enumerate}
		\item
		\[
		\|\widehat f\|_\infty\leq \|f\|_{L_1};
		\]
		
		\item \(\widehat f\) непрерывна;
		
		\item
		\[
		\widehat f(y)\to0,
		\qquad |y|\to\infty.
		\]
	\end{enumerate}
	
	Таким образом,
	\[
	\mathcal F:L_1(\mathbb R)\to C_0(\mathbb R)
	\]
	--- линейный ограниченный оператор.
\end{theorem}

\begin{proof}
	\begin{enumerate}
		\item
		\textbf{Хотим оценить норму \(\widehat f\).}
		
		Для любого \(y\in\mathbb R\)
		\[
		|\widehat f(y)|
		=
		\left|
		\int_{\mathbb R}f(x)e^{-ixy}\,dx
		\right|
		\leq
		\int_{\mathbb R}|f(x)|\,dx
		=
		\|f\|_{L_1}.
		\]
		Следовательно,
		\[
		\|\widehat f\|_\infty\leq\|f\|_{L_1}.
		\]
		
		\item
		\textbf{Хотим доказать непрерывность \(\widehat f\).}
		
		Пусть \(y_n\to y\). Тогда
		\[
		f(x)e^{-ixy_n}\to f(x)e^{-ixy}
		\]
		для почти всех \(x\), причем
		\[
		|f(x)e^{-ixy_n}|=|f(x)|\in L_1.
		\]
		По теореме Лебега
		\[
		\widehat f(y_n)\to\widehat f(y).
		\]
		
		\item
		\textbf{Хотим показать убывание на бесконечности.}
		
		Для \(\varphi\in C_c^1(\mathbb R)\) интегрированием по частям получаем
		\[
		\widehat\varphi(y)
		=
		\frac{1}{iy}
		\int_{\mathbb R}\varphi'(x)e^{-ixy}\,dx,
		\]
		поэтому
		\[
		|\widehat\varphi(y)|
		\leq
		\frac{\|\varphi'\|_{L_1}}{|y|}
		\to0.
		\]
		
		Теперь для \(f\in L_1\) выбираем
		\[
		\varphi\in C_c^1(\mathbb R),
		\qquad
		\|f-\varphi\|_{L_1}<\varepsilon.
		\]
		Тогда
		\[
		|\widehat f(y)|
		\leq
		|\widehat{f-\varphi}(y)|+|\widehat\varphi(y)|
		\leq
		\varepsilon+|\widehat\varphi(y)|.
		\]
		При достаточно больших \(|y|\)
		\[
		|\widehat\varphi(y)|<\varepsilon,
		\]
		следовательно,
		\[
		|\widehat f(y)|<2\varepsilon.
		\]
	\end{enumerate}
\end{proof}

\begin{corollary}
	Если
	\[
	f_n\to f
	\quad\text{в }L_1,
	\]
	то
	\[
	\widehat f_n\to\widehat f
	\quad\text{равномерно на }\mathbb R.
	\]
\end{corollary}

\begin{proof}
	Действительно,
	\[
	\|\widehat f_n-\widehat f\|_\infty
	=
	\|\widehat{f_n-f}\|_\infty
	\leq
	\|f_n-f\|_{L_1}
	\to0.
	\]
\end{proof}

\begin{theorem}[Формула умножения]
	Если
	\[
	f,g\in L_1(\mathbb R),
	\]
	то
	\[
	\int_{\mathbb R}\widehat f(y)g(y)\,dy
	=
	\int_{\mathbb R}f(x)\widehat g(x)\,dx.
	\]
\end{theorem}

\begin{proof}
	\textbf{Хотим поменять порядок интегрирования.}
	Для этого проверим абсолютную интегрируемость:
	\[
	\int_{\mathbb R}\int_{\mathbb R}
	|f(x)g(y)e^{-ixy}|\,dx\,dy
	=
	\|f\|_{L_1}\|g\|_{L_1}
	<\infty.
	\]
	По теореме Фубини
	\[
	\begin{aligned}
		\int_{\mathbb R}\widehat f(y)g(y)\,dy
		&=
		\int_{\mathbb R}
		\left(
		\int_{\mathbb R}f(x)e^{-ixy}\,dx
		\right)g(y)\,dy
		\\
		&=
		\int_{\mathbb R}f(x)
		\left(
		\int_{\mathbb R}g(y)e^{-ixy}\,dy
		\right)dx
		\\
		&=
		\int_{\mathbb R}f(x)\widehat g(x)\,dx.
	\end{aligned}
	\]
\end{proof}

\begin{definition}[Свертка]
	Пусть
	\[
	f,g\in L_1(\mathbb R).
	\]
	Их сверткой называется функция
	\[
	(f*g)(x)
	=
	\int_{\mathbb R}f(y)g(x-y)\,dy.
	\]
\end{definition}

\begin{theorem}
	Если
	\[
	f,g\in L_1(\mathbb R),
	\]
	то
	\[
	f*g\in L_1(\mathbb R)
	\]
	и
	\[
	\|f*g\|_{L_1}
	\leq
	\|f\|_{L_1}\|g\|_{L_1}.
	\]
\end{theorem}

\begin{proof}
	\textbf{Хотим показать интегрируемость свертки.}
	
	По теореме Тонелли
	\[
	\begin{aligned}
		\|f*g\|_{L_1}
		&\leq
		\int_{\mathbb R}\int_{\mathbb R}
		|f(y)||g(x-y)|\,dy\,dx
		\\
		&=
		\int_{\mathbb R}|f(y)|
		\left(
		\int_{\mathbb R}|g(x-y)|\,dx
		\right)dy
		\\
		&=
		\|f\|_{L_1}\|g\|_{L_1}.
	\end{aligned}
	\]
\end{proof}

\begin{theorem}[Преобразование Фурье свертки]
	Если
	\[
	f,g\in L_1(\mathbb R),
	\]
	то
	\[
	\widehat{f*g}
	=
	\widehat f\,\widehat g.
	\]
\end{theorem}

\begin{proof}
	\textbf{Хотим раскрыть свертку и поменять порядок интегрирования.}
	
	Так как
	\[
	\int_{\mathbb R}\int_{\mathbb R}
	|f(y)g(x-y)|\,dy\,dx
	=
	\|f\|_{L_1}\|g\|_{L_1}
	<\infty,
	\]
	по теореме Фубини
	\[
	\begin{aligned}
		\widehat{f*g}(z)
		&=
		\int_{\mathbb R}
		\left(
		\int_{\mathbb R}f(y)g(x-y)\,dy
		\right)e^{-ixz}\,dx
		\\
		&=
		\int_{\mathbb R}f(y)
		\left(
		\int_{\mathbb R}g(x-y)e^{-ixz}\,dx
		\right)dy.
	\end{aligned}
	\]
	
	\textbf{Хотим выделить преобразование Фурье функции \(g\).}
	Делаем замену
	\[
	t=x-y.
	\]
	Тогда
	\[
	\begin{aligned}
		\widehat{f*g}(z)
		&=
		\int_{\mathbb R}f(y)
		\left(
		\int_{\mathbb R}g(t)e^{-i(t+y)z}\,dt
		\right)dy
		\\
		&=
		\left(
		\int_{\mathbb R}f(y)e^{-iyz}\,dy
		\right)
		\left(
		\int_{\mathbb R}g(t)e^{-itz}\,dt
		\right)
		\\
		&=
		\widehat f(z)\widehat g(z).
	\end{aligned}
	\]
\end{proof}

\subsection{Преобразование Фурье в пространстве \texorpdfstring{\(L_2(\mathbb R)\)}{L2(R)}}

Рассмотрим пространство Шварца \(S(\mathbb R)\).
Используем без доказательства следующие факты:
\[
\overline{S(\mathbb R)}^{\,L_2}=L_2(\mathbb R),
\qquad
\mathcal F:S(\mathbb R)\to S(\mathbb R)
\]
--- биекция, причем
\[
\|\widehat f\|_{L_2}=\|f\|_{L_2},
\qquad f\in S(\mathbb R).
\]

Здесь
\[
\widehat f(y)
=
\frac{1}{\sqrt{2\pi}}
\int_{\mathbb R}f(x)e^{-ixy}\,dx.
\]

Так как \(S(\mathbb R)\) плотно в \(L_2(\mathbb R)\),
а \(\mathcal F\) является изометрией на \(S\),
она единственным образом продолжается до линейной изометрии
\[
\mathcal F:L_2(\mathbb R)\to L_2(\mathbb R).
\]

\begin{theorem}[Сохранение скалярного произведения]
	Для любых
	\[
	f,g\in L_2(\mathbb R)
	\]
	выполнено
	\[
	(f,g)=(\widehat f,\widehat g).
	\]
\end{theorem}

\begin{proof}
	\textbf{Сначала докажем равенство для \(f,g\in S\).}
	
	По формуле обращения
	\[
	\overline{g(x)}
	=
	\frac{1}{\sqrt{2\pi}}
	\int_{\mathbb R}
	\overline{\widehat g(y)}e^{-ixy}\,dy.
	\]
	
	По теореме Фубини
	\[
	\begin{aligned}
		(f,g)
		&=
		\int_{\mathbb R}f(x)\overline{g(x)}\,dx
		\\
		&=
		\frac{1}{\sqrt{2\pi}}
		\int_{\mathbb R}\int_{\mathbb R}
		f(x)\overline{\widehat g(y)}e^{-ixy}\,dy\,dx
		\\
		&=
		\int_{\mathbb R}
		\widehat f(y)\overline{\widehat g(y)}\,dy
		\\
		&=
		(\widehat f,\widehat g).
	\end{aligned}
	\]
	
	\textbf{Теперь переходим к \(f,g\in L_2\).}
	
	Берем
	\[
	f_n,g_n\in S,
	\qquad
	f_n\to f,
	\qquad
	g_n\to g
	\quad\text{в }L_2.
	\]
	
	Из непрерывности \(\mathcal F\)
	\[
	\widehat f_n\to\widehat f,
	\qquad
	\widehat g_n\to\widehat g.
	\]
	
	По непрерывности скалярного произведения
	\[
	(f_n,g_n)\to(f,g),
	\qquad
	(\widehat f_n,\widehat g_n)
	\to
	(\widehat f,\widehat g).
	\]
	
	Так как
	\[
	(f_n,g_n)
	=
	(\widehat f_n,\widehat g_n),
	\]
	переходом к пределу получаем
	\[
	(f,g)=(\widehat f,\widehat g).
	\]
\end{proof}

\begin{theorem}[Основные свойства преобразования Фурье в \(L_2\)]
	Для
	\[
	\mathcal F:L_2(\mathbb R)\to L_2(\mathbb R)
	\]
	выполнено:
	\begin{enumerate}
		\item
		\[
		\operatorname{Im}\mathcal F=L_2(\mathbb R);
		\]
		
		\item
		\[
		\ker\mathcal F=\{0\};
		\]
		
		\item \(\mathcal F\) не является компактным оператором;
		
		\item
		\[
		\mathcal F^4=I.
		\]
	\end{enumerate}
\end{theorem}

\begin{proof}
	\begin{enumerate}
		\item
		\textbf{Хотим доказать сюръективность.}
		
		Пусть
		\[
		f\in L_2.
		\]
		Так как \(S\) плотно в \(L_2\), существуют
		\[
		f_n\in S,
		\qquad
		f_n\to f.
		\]
		
		Так как
		\[
		\mathcal F:S\to S
		\]
		--- биекция, для каждого \(n\) существует
		\[
		g_n\in S:
		\qquad
		\widehat g_n=f_n.
		\]
		
		По изометричности
		\[
		\|g_n-g_m\|_{L_2}
		=
		\|\widehat g_n-\widehat g_m\|_{L_2}
		=
		\|f_n-f_m\|_{L_2}.
		\]
		
		Поэтому \(\{g_n\}\) фундаментальна, а значит
		\[
		g_n\to g
		\]
		для некоторого \(g\in L_2\).
		
		По непрерывности \(\mathcal F\)
		\[
		\widehat g_n\to\widehat g.
		\]
		Но
		\[
		\widehat g_n=f_n\to f,
		\]
		следовательно,
		\[
		\widehat g=f.
		\]
		
		Значит,
		\[
		\operatorname{Im}\mathcal F=L_2.
		\]
		
		\item
		\textbf{Хотим доказать инъективность.}
		
		Если
		\[
		\widehat f=0,
		\]
		то по изометричности
		\[
		\|f\|_{L_2}
		=
		\|\widehat f\|_{L_2}
		=
		0.
		\]
		Следовательно,
		\[
		f=0,
		\]
		то есть
		\[
		\ker\mathcal F=\{0\}.
		\]
		
		\item
		\textbf{Хотим доказать, что \(\mathcal F\) не компактен.}
		
		Предположим противное: \(\mathcal F\) компактен.
		Тогда композиция
		\[
		\mathcal F^4
		\]
		также компактна.
		
		Но
		\[
		\mathcal F^4=I.
		\]
		
		Получаем, что тождественный оператор
		\[
		I:L_2\to L_2
		\]
		компактен, что невозможно, поскольку \(L_2(\mathbb R)\)
		бесконечномерно.
		
		Следовательно, \(\mathcal F\) не компактен.
		
		\item
		\textbf{Хотим вычислить четвертую степень \(\mathcal F\).}
		
		На пространстве \(S\) из формулы обращения следует
		\[
		\mathcal F^2f(x)=f(-x).
		\]
		
		Поэтому
		\[
		\mathcal F^4f=f,
		\qquad f\in S.
		\]
		
		Так как \(S\) плотно в \(L_2\), а \(\mathcal F\) непрерывен,
		равенство продолжается на всё \(L_2\):
		\[
		\mathcal F^4=I.
		\]
	\end{enumerate}
\end{proof}

\begin{idea}
	\textbf{Зачем вообще нужно пространство Шварца и что мы здесь делаем.}
	
	В пространстве \(L_1(\mathbb R)\) преобразование Фурье можно определить напрямую:
	\[
	\widehat f(y)
	=
	\int_{\mathbb R}f(x)e^{-ixy}\,dx,
	\]
	поскольку
	\[
	\int_{\mathbb R}|f(x)e^{-ixy}|\,dx
	=
	\int_{\mathbb R}|f(x)|\,dx
	<\infty.
	\]
	
	\textbf{В \(L_2\) так сделать нельзя.}
	
	Из
	\[
	f\in L_2(\mathbb R)
	\]
	вообще не следует, что
	\[
	f\in L_1(\mathbb R).
	\]
	Поэтому для произвольного \(f\in L_2\) интеграл
	\[
	\int_{\mathbb R}f(x)e^{-ixy}\,dx
	\]
	может не существовать как обычный интеграл.
	
	То есть наша цель:
	\[
	\boxed{\text{определить преобразование Фурье для всех }f\in L_2}
	\]
	но напрямую по интегральной формуле это сделать нельзя.
	
	\medskip
	
	\textbf{1. Поэтому сначала берем ``хорошие'' функции.}
	
	Рассматриваем пространство Шварца
	\[
	S(\mathbb R)
	=
	\left\{
	f\in C^\infty(\mathbb R):
	\sup_{x\in\mathbb R}|x^m f^{(n)}(x)|<\infty
	\quad
	\forall m,n\geq0
	\right\}.
	\]
	
	Смысл условия: функции из \(S\) гладкие и вместе со всеми производными
	убывают на бесконечности быстрее любой степени.
	
	В частности,
	\[
	S(\mathbb R)\subset L_1(\mathbb R)\cap L_2(\mathbb R).
	\]
	
	Поэтому для \(f\in S\) интегральная формула преобразования Фурье
	точно имеет смысл:
	\[
	\widehat f(y)
	=
	\frac{1}{\sqrt{2\pi}}
	\int_{\mathbb R}f(x)e^{-ixy}\,dx.
	\]
	
	То есть сначала мы умеем корректно работать с
	\[
	\mathcal F:S\to S.
	\]
	
	\medskip
	
	\textbf{2. Почему именно \(S\), а не какое-то случайное подпространство?}
	
	У пространства Шварца есть два нужных нам свойства,
	которые используем без доказательства:
	\[
	\boxed{\overline{S}^{\,L_2}=L_2}
	\]
	и
	\[
	\boxed{\mathcal F:S\to S\text{ --- биекция}.}
	\]
	
	Первое означает, что любую функцию
	\[
	f\in L_2
	\]
	можно приблизить функциями из \(S\):
	\[
	f_n\in S,
	\qquad
	f_n\to f
	\quad\text{в }L_2.
	\]
	
	Поэтому, если мы хорошо определим преобразование Фурье на \(S\)
	и докажем его непрерывность в норме \(L_2\),
	то сможем распространить его на всё \(L_2\).
	
	\medskip
	
	\textbf{3. Главное свойство на \(S\) --- Фурье сохраняет норму.}
	
	Для \(f\in S\) доказывается равенство Планшереля:
	\[
	\|\widehat f\|_{L_2}
	=
	\|f\|_{L_2}.
	\]
	
	То есть
	\[
	\mathcal F:S\to S
	\]
	является изометрией.
	
	Это именно то, что позволяет перейти от \(S\) ко всему \(L_2\).
	
	\medskip
	
	\textbf{4. Как по \(S\) определить Фурье произвольной функции \(f\in L_2\)?}
	
	Берем последовательность хороших функций
	\[
	f_n\in S,
	\qquad
	f_n\to f
	\quad\text{в }L_2.
	\]
	
	У каждой \(f_n\) преобразование Фурье уже определено обычным интегралом:
	\[
	\widehat f_n
	=
	\mathcal Ff_n.
	\]
	
	Хотим определить
	\[
	\widehat f
	\]
	как предел этих преобразований.
	
	Из изометричности
	\[
	\|\widehat f_n-\widehat f_m\|_{L_2}
	=
	\|f_n-f_m\|_{L_2}.
	\]
	
	Так как
	\[
	f_n\to f,
	\]
	последовательность \(\{f_n\}\) фундаментальна.
	Следовательно, фундаментальна и
	\[
	\{\widehat f_n\}.
	\]
	
	Пространство \(L_2\) полно, поэтому существует
	\[
	g\in L_2:
	\qquad
	\widehat f_n\to g.
	\]
	
	Именно этот предел и объявляем преобразованием Фурье функции \(f\):
	\[
	\boxed{
		\mathcal Ff
		:=
		\lim_{n\to\infty}\mathcal Ff_n
	}.
	\]
	
	То есть для общей функции \(f\in L_2\)
	преобразование Фурье определяется уже
	\textbf{не обязательно обычным интегралом},
	а пределом в норме \(L_2\).
	
	\medskip
	
	\textbf{5. Почему результат не зависит от того, как именно мы приблизили \(f\)?}
	
	Пусть есть две последовательности
	\[
	f_n\to f,
	\qquad
	g_n\to f,
	\qquad
	f_n,g_n\in S.
	\]
	
	Тогда
	\[
	\|f_n-g_n\|_{L_2}\to0.
	\]
	
	По изометричности
	\[
	\|\widehat f_n-\widehat g_n\|_{L_2}
	=
	\|f_n-g_n\|_{L_2}
	\to0.
	\]
	
	Значит, обе последовательности преобразований дают один и тот же предел.
	
	Следовательно, \(\mathcal Ff\) определено корректно.
	
	\medskip
	
	\textbf{Итого вся конструкция выглядит так:}
	\[
	\boxed{
		S
		\overset{\text{плотно}}{\subset}
		L_2
	}
	\]
	\[
	\boxed{
		\mathcal F:S\to S
		\text{ уже определено обычным интегралом}
	}
	\]
	\[
	\boxed{
		\|\mathcal Ff\|_2=\|f\|_2
	}
	\]
	\[
	\Downarrow
	\]
	\[
	\boxed{
		\mathcal F
		\text{ единственным образом продолжается }
		S\to S
		\quad\text{до}\quad
		L_2\to L_2.
	}
	\]
	
	После этого мы уже изучаем свойства полученного оператора на \(L_2\):
	\[
	(f,g)=(\widehat f,\widehat g),
	\qquad
	\ker\mathcal F=\{0\},
	\qquad
	\operatorname{Im}\mathcal F=L_2,
	\qquad
	\mathcal F^4=I,
	\]
	и отсюда, в частности, следует, что \(\mathcal F\) не компактен.
	
	Таким образом, пространство Шварца здесь нужно не само по себе:
	это просто удобное плотное множество ``хороших'' функций,
	на котором преобразование Фурье сначала можно определить
	и изучить обычными интегралами, а затем по непрерывности
	перенести на всё \(L_2\).
\end{idea}\newpage
	\hypertarget{bilet20}{}\setupnav{19}{} % --- Содержимое файла materials/14_2.tex ---

\subsection{Операторы Гильберта--Шмидта.}

\begin{remark}
	TODO: сюда дописывается содержимое билета в том же стиле, что и в конспекте по функциональному анализу--1.
	Сначала фиксируются нужные определения и формулировки теорем, затем подробно расписываются доказательства без пропуска логических переходов.
\end{remark}
\newpage
	
\end{document}